% bismi ALLAH ARRAHMAN ARRAHIM

















































































































\begin{theorem}[Loi effective de courbure locale]
\label{thm:effective-law}
Soit $u^\star \in \mathbb R^d$ un point fixe local du noyau PCH, et supposons qu'au voisinage de $u^\star$ la fonctionnelle effective puisse être écrite, à l'ordre quadratique dominant, sous la forme
\[
\Phi_{\mathrm{eff}}(u)
=
\frac{\mu_{\mathrm{eff}}}{2}\,\|u-u^\star\|^2
+
o\!\left(\|u-u^\star\|^2\right),
\qquad \mu_{\mathrm{eff}}>0.
\]
Pour $\eta>0$, considérons l'opérateur proximal associé
\[
T_\eta(x)
=
\arg\min_u
\left\{
\Phi_{\mathrm{eff}}(u)
+
\frac{1}{2\eta}\|u-x\|^2
\right\}.
\]
Alors, au voisinage de $u^\star$,
\[
T_\eta(x)-u^\star
=
\frac{1}{1+\eta \mu_{\mathrm{eff}}}\,(x-u^\star)
+
o\!\left(\|x-u^\star\|\right),
\]
et le facteur local de contraction vaut donc
\[
\rho(\eta)=\frac{1}{1+\eta\mu_{\mathrm{eff}}}.
\]
\end{theorem}

\begin{proof}
On écrit l'optimalité de $T_\eta(x)$. En notant $u=T_\eta(x)$, on a
\[
0 \in \partial \Phi_{\mathrm{eff}}(u)+\frac{1}{\eta}(u-x).
\]
Sous l'hypothèse quadratique locale,
\[
\nabla \Phi_{\mathrm{eff}}(u)=\mu_{\mathrm{eff}}(u-u^\star)+o(\|u-u^\star\|).
\]
Ainsi,
\[
\mu_{\mathrm{eff}}(u-u^\star)+\frac{1}{\eta}(u-x)
=
o(\|u-u^\star\|).
\]
En réarrangeant,
\[
\left(\mu_{\mathrm{eff}}+\frac{1}{\eta}\right)(u-u^\star)
=
\frac{1}{\eta}(x-u^\star)
+
o(\|u-u^\star\|).
\]
Comme $\mu_{\mathrm{eff}}+\frac{1}{\eta}>0$, on obtient
\[
u-u^\star
=
\frac{1}{1+\eta\mu_{\mathrm{eff}}}(x-u^\star)
+
o\!\left(\|x-u^\star\|\right),
\]
ce qui prouve la formule annoncée.
\end{proof}

\begin{corollary}[Identification de la constante observée]
\label{cor:constant-identification}
Si l'on observe exactement
\[
\rho(1)=\frac{5}{22},
\]
alors la courbure effective est nécessairement
\[
\mu_{\mathrm{eff}}
=
\frac{1-\rho(1)}{\rho(1)}
=
\frac{17}{5}.
\]
Par conséquent,
\[
\rho(2)
=
\frac{1}{1+2\mu_{\mathrm{eff}}}
=
\frac{1}{1+\frac{34}{5}}
=
\frac{5}{39}.
\]
Réciproquement, si $\rho(2)=5/39$, alors la même valeur $\mu_{\mathrm{eff}}=17/5$ est imposée.
\end{corollary}

\begin{proof}
La première identité suit de
\[
\rho(\eta)=\frac{1}{1+\eta\mu_{\mathrm{eff}}}
\quad\Longleftrightarrow\quad
\mu_{\mathrm{eff}}
=
\frac{1-\rho(\eta)}{\eta\,\rho(\eta)}.
\]
En posant $\eta=1$ et $\rho(1)=5/22$, on obtient
\[
\mu_{\mathrm{eff}}
=
\frac{1-\frac{5}{22}}{\frac{5}{22}}
=
\frac{\frac{17}{22}}{\frac{5}{22}}
=
\frac{17}{5}.
\]
Puis, pour $\eta=2$,
\[
\rho(2)
=
\frac{1}{1+2\cdot \frac{17}{5}}
=
\frac{1}{\frac{39}{5}}
=
\frac{5}{39}.
\]
La réciproque est identique en sens inverse.
\end{proof}

\begin{proposition}[Loi fermée de la dynamique locale]
\label{prop:closed-law}
Sous les hypothèses du théorème \ref{thm:effective-law}, si les diagnostics expérimentaux sont cohérents avec
\[
\rho(1)=\frac{5}{22}
\qquad\text{et}\qquad
\rho(2)=\frac{5}{39},
\]
alors la loi locale effective est nécessairement
\[
\rho(\eta)=\frac{1}{1+\frac{17}{5}\eta}
=
\frac{5}{5+17\eta}.
\]
\end{proposition}

\begin{proof}
D'après le théorème \ref{thm:effective-law}, la contraction locale a la forme
\[
\rho(\eta)=\frac{1}{1+\eta\mu_{\mathrm{eff}}}.
\]
Le corollaire \ref{cor:constant-identification} identifie alors
\[
\mu_{\mathrm{eff}}=\frac{17}{5}.
\]
En substituant,
\[
\rho(\eta)=\frac{1}{1+\frac{17}{5}\eta}
=
\frac{5}{5+17\eta}.
\]
\end{proof}

\begin{remark}[Portée exacte de ce résultat]
Ce résultat prouve une \emph{loi effective} conditionnelle : si la dynamique PCH est localement bien approchée par un prox quadratique isotrope, alors la constante observée n'est pas accidentelle mais impose une courbure effective unique $\mu_{\mathrm{eff}}=17/5$.

En revanche, ce n'est pas encore une preuve que les données expérimentales satisfont exactement cette hypothèse. Pour cela, il faudrait démontrer analytiquement ou valider numériquement que le Jacobien local du noyau PCH est bien de la forme
\[
D T_\eta(u^\star)=\frac{1}{1+\eta\mu_{\mathrm{eff}}}I
\]
sur le sous-espace actif pertinent, ou au moins établir une approximation spectrale suffisamment précise.
\end{remark}




%%%%%%%%%%%%%%%%%%%%%%%%%%%%%%%%%%%%%%%%%%%%%%%%%%%%%%%%%%%%%%%%%%%%%%%%%%%%%%%%%%%%%%%%%%%%%

\subsection{Vérification expérimentale de la loi locale}
On fixe le contexte \(c\), la graine, la dimension latente, la taille de calibration et la règle de routage.
Pour chaque \(\eta\in\{1,2\}\), on calcule :
\[
u^\star=\Ppch(w;c),\qquad
\widehat\rho_{\mathrm{lin}}(\eta)=\|DT_\eta(u^\star;w,c)\|_2,
\qquad
\widehat\rho_{\mathrm{use}}(\eta)=\frac{1}{K}\sum_{r=1}^K
\frac{\|u^{(r+1)}-u^\star\|}{\|u^{(r)}-u^\star\|},
\]
en ne conservant que les premières itérations utiles avant le plancher numérique.
On estime ensuite
\[
\widehat\mu(\eta)=\frac{1}{\eta}\left(\frac{1}{\widehat\rho_{\mathrm{lin}}(\eta)}-1\right).
\]
La loi est acceptée si \(\widehat\mu(1)\approx \widehat\mu(2)\) et si
\[
\widehat\rho_{\mathrm{lin}}(\eta)\approx \frac{1}{1+\eta\widehat\mu}
\]
pour tous les régimes testés, avec zéro violation de décroissance d’énergie et des résidus de point fixe au voisinage de la précision machine.

%%%%%%%%%%%%%%%%%%%%%%%%%%%%%%%%%%%%%%%%%%%%%%%%%%%%%%%%%%%%%%%%%%%%%%%%%%%%%%%%%%%%%%%%%%%






\subsection{Experience "verify_pch_local_law" : Description rigoureuse de l'expérience, du pipeline et des résultats}

\paragraph{Objet de l'expérience.}
L'expérience étudie le noyau PCH appliqué à des embeddings textuels réels, dans un cadre entièrement contrôlé : encodeur figé, réduction PCA, ensembles contextuels admissibles construits comme enveloppes convexes de calibrations, et résolution exacte des sous-problèmes convexes par QP. Le pipeline expérimental implémente explicitement
\[
w \xrightarrow[\text{SentenceTransformer}]{E_{\theta_0}}
e \xrightarrow[\text{PCA}]{\text{latent}}
z \xrightarrow[\text{hull du contexte } c]{\mathcal M_c}
\Ppch(w;c)
\quad\text{puis}\quad
u^{(r+1)}=\Tpch(u^{(r)};w,c),
\]
avec un diagnostic de contraction effectué uniquement sur les premières itérations utiles, avant domination du bruit numérique, et une linéarisation locale estimée par jacobien à différences finies au voisinage du point fixe.

\paragraph{Protocole.}
Les expériences utilisent le corpus \texttt{20 Newsgroups} dans un sous-ensemble à quatre classes pour le régime principal, et une extension à vingt classes pour tester la difficulté du routage. L’encodeur est figé (\texttt{sentence-transformers/all-MiniLM-L6-v2}), la PCA est déterministe, les contextes sont définis par
\[
\mathcal M_c=\operatorname{conv}\{E_{\theta_0}(x):x\in D_c\},
\]
et l’énergie testée est
\[
J_{\mathrm{PCH}}(u;w,c)=\frac12\|u-w\|^2+\Phi_{\mathrm{PCH},c}(u),
\qquad
\Phi_{\mathrm{PCH},c}(u)=\lambda_{\mathrm{ctx}}\operatorname{dist}(u,\mathcal M_c)^2+\lambda_{\mathrm{ridge}}\|u\|^2,
\]
avec, dans la version finale du script, \(\lambda_{\mathrm{ridge}}=\lambda_{\mathrm{ch}}+\lambda_{\mathrm{reg}}=1.2\). Les runs conservés incluent notamment \(\eta\in\{1,2\}\), \(d\in\{16,32,48,63,64,80,128\}\), \(n_{\mathrm{calib}}\in\{16,32,48,64,80,128\}\), et plusieurs valeurs de \(\lambda_{\mathrm{ctx}}\). Le calcul du prox et du minimiseur est exact au sens d’un solveur convexe, sans descente de gradient ni heuristique cachée.

\paragraph{Résultats expérimentaux consolidés.}
Les runs consolidés montrent, de manière récurrente, trois faits robustes. Premièrement, la décroissance d’énergie n’est jamais violée : tous les rapports conservés affichent \(\texttt{all\_energy\_monotonic\_violations}=0\). Deuxièmement, les résidus de point fixe sont au voisinage de la précision machine, typiquement entre \(10^{-17}\) et \(10^{-16}\), avec des maxima plus élevés dans quelques réglages plus exigeants, sans perte de stabilité du schéma. Troisièmement, la norme de la linéarisation locale du prox est pratiquement constante d’un run à l’autre. Dans le régime \(\eta=1\), elle se situe autour de
\[
\rho_{\mathrm{lin}}\approx 0.227272727\dots,
\]
et dans le régime \(\eta=2\), le protocole théorique et les vérifications associées indiquent la valeur
\[
\rho_{\mathrm{lin}}\approx 0.128205128\dots,
\]
ce qui confirme la baisse attendue du facteur de contraction quand le paramètre proximal augmente. Cette stabilité est corroborée par les diagnostics de \(\mu\)-consistance, qui donnent \(\mu\simeq 3.4\), compatible avec la formule
\[
\rho(\eta)=\frac{1}{1+\eta\mu},
\qquad
\mu=1+2(\lambda_{\mathrm{ch}}+\lambda_{\mathrm{reg}})=3.4.
\]

Les runs représentatifs montrent aussi que la purification rapproche systématiquement les points de leur hull admissible. Par exemple, avec \(d=32\), \(n_{\mathrm{calib}}=80\), \(\lambda_{\mathrm{ctx}}=0.2\), le routage par distance de projection donne une réduction nette de distance et un gain nul ou faible sur le probe selon le régime, tout en conservant la stabilité numérique totale. Dans un régime plus favorable, \(d=128\), \(n_{\mathrm{calib}}=128\), \(\lambda_{\mathrm{ctx}}=1.7\), \(\texttt{max\_iters}=5\), le passage de l’embedding brut à l’embedding PCH améliore l’accuracy de \(0.640625\) à \(0.7109375\), avec une réduction de distance très forte au hull et zéro violation d’énergie. Les contextes 4-classes à plus faible calibration peuvent au contraire produire un gain nul ou négatif, malgré une réduction géométrique claire, ce qui montre que l’effet géométrique du noyau n’implique pas automatiquement un gain aval monotone. Les résultats extraits du fichier \texttt{pch\_all\_exps.txt} confirment ce comportement contrasté, tout en conservant la stabilité du noyau et la monotonie énergétique. 

\paragraph{Interprétation.}
L’interprétation correcte est la suivante : le noyau PCH agit comme un opérateur proximal exact, stable et fortement contractif au voisinage du point fixe. Il réalise une purification géométrique contextuelle mesurable par la baisse de distance au hull, mais il ne crée pas à lui seul une séparabilité universelle. La partie la plus solide du résultat est la loi locale de contraction : la linéarisation observée est cohérente, à la précision machine, avec la borne théorique
\[
\rho_{\mathrm{lin}}\le \frac{1}{1+\eta\mu},
\]
et l’expérience confirme numériquement cette loi pour \(\eta=1\) et \(\eta=2\). En revanche, l’égalité entre “useful ratios” et norme de linéarisation n’est pas une loi exacte : les useful ratios sont des diagnostics pré-asymptotiques, systématiquement un peu inférieurs à la linéarisation locale, avec un écart dépendant du régime. Ce point est important, car il explique pourquoi certaines vérifications automatisées échouent alors que la dynamique de fond reste parfaitement cohérente. 

\paragraph{Conséquences.}
La conséquence théorique principale est que la partie conditionnelle du noyau PCH est validée expérimentalement : existence numérique du minimiseur, unicité effective à la précision machine, contraction locale, convergence du prox, stabilité numérique, et réduction de distance au contexte admissible. En revanche, la théorie du choix automatique de contexte n’est pas démontrée par ces expériences. Le routage par énergie PCH apparaît trop bruité pour être un routeur universel fiable dans les régimes testés, tandis que le routage par distance de projection est empiriquement plus robuste. Il faut donc distinguer très nettement :
\[
\text{PCH conditionnel à }c
\quad\text{(validé)}
\qquad\text{et}\qquad
\text{routage universel}
\quad\text{(non validé).}
\]

\paragraph{Verdict sur la loi confirmée.}
La loi confirmée est la loi de contraction locale du prox, avec dépendance explicite en \(\eta\) :
\[
\rho(\eta)=\frac{1}{1+\eta\mu},
\qquad
\mu=3.4,
\]
soit numériquement
\[
\rho(1)=\frac{1}{4.4}=\frac{5}{22}\approx 0.227272727,
\qquad
\rho(2)=\frac{1}{7.8}=\frac{5}{39}\approx 0.128205128.
\]
Cette loi est soutenue par les diagnostics de linéarisation, par la \(\mu\)-consistance du script, par l’absence totale de violation énergétique, et par les résidus de point fixe au voisinage de \(10^{-17}\) à \(10^{-16}\). En revanche, la loi “les useful ratios doivent coïncider exactement avec la norme linéarisée” n’est pas confirmée ; les useful ratios doivent être lus comme un estimateur pré-asymptotique, pas comme une identité analytique.




%%%%%%%%%%%%%%%%%%%%%%%%%%%%%%%%%%%%%%%%%%%%%%%%%%%%%%%%%%%%%%%%%%%%%%%%%%%%%%%%%%%%%%%%%%%%%

%%%%%%%%%%%EXPERIENCE / pch_cas3_sweep.py

\section{Description rigoureuse des expériences, du pipeline, des protocoles et des résultats}

\subsection{Objet général}

On étudie le noyau PCH sur des embeddings textuels réels, dans deux régimes expérimentaux complémentaires :

\begin{itemize}
    \item \textbf{ex1 :} corpus \texttt{20news20} (20 classes), encodeur figé \texttt{sentence-transformers/all-MiniLM-L6-v2}, PCA de dimension effective 48, 48 exemples de calibration par contexte, 32 exemples de probe train/test par contexte, \(\eta=1\), \(\lambda_{\mathrm{ctx}}=1.7\), \(\lambda_{\mathrm{ch}}=0.4\), \(\lambda_{\mathrm{reg}}=0.8\), soit \(\lambda_{\mathrm{ridge}}=1.2\).
    \item \textbf{ex2 :} sweep sur \texttt{20news4} (4 classes), même encodeur ; PCA ajustée une seule fois sur les embeddings du sous-ensemble de calibration du split d’apprentissage, puis réutilisée sans réentraînement ; 48 calibrations par contexte, 32 exemples de probe train/test par contexte, \(n_{\text{train}}=n_{\text{test}}=128\) pour les splits de probe observés dans les logs, \(\eta\in\{1,2\}\), \(\lambda_{\mathrm{ctx}}=0.2\), \(\lambda_{\mathrm{ch}}=0\), \(\lambda_{\mathrm{reg}}\in\{0.2,0.8\}\), avec régularisation \texttt{ridge} ou \texttt{none}.
\end{itemize}

Le cœur de l'expérience consiste à vérifier simultanément :
\begin{enumerate}
    \item la convergence proximale exacte ;
    \item la décroissance énergétique ;
    \item la stabilité du point fixe ;
    \item la loi locale de contraction ;
    \item la réduction géométrique vis-à-vis du hull admissible ;
    \item la qualité du routage de contexte ;
    \item l'effet aval sur un probe linéaire.
\end{enumerate}

\subsection{Pipeline expérimental}

Pour chaque texte \(x\), on applique successivement :

\[
x \xrightarrow{E_{\theta_0}} e \xrightarrow{\mathrm{PCA}} z \xrightarrow{\mathcal M_c=\operatorname{conv}(D_c)} \Ppch(w;c)
\quad\text{puis}\quad
u^{(r+1)}=\Tpch(u^{(r)};w,c),
\]
où :

\begin{itemize}
    \item \(E_{\theta_0}\) est un encodeur figé de type SentenceTransformer ;
    \item la PCA est déterministe ;
    \item \(D_c\) est l’ensemble des embeddings de calibration appartenant au contexte \(c\) ; dans les expériences rapportées ici, \(D_c\) contient 48 points par contexte ;
    \item \(\mathcal M_c = \operatorname{conv}(D_c)\) est l’enveloppe convexe de ces embeddings de calibration dans l’espace latent ;
    \item \(\Ppch(w;c)\) est le minimiseur exact du problème PCH ;
    \item \(\Tpch\) est l’itération proximale exacte ;
    \item la contraction est mesurée uniquement sur les premières itérations utiles, avant domination du bruit machine ;
    \item la linéarisation locale est estimée par jacobien à différences finies au voisinage du point fixe.
\end{itemize}

Le solveur utilisé est un QP convexe exact via \texttt{cvxpy + OSQP}. Aucune descente de gradient ni approximation stochastique n’intervient dans le calcul du prox.

\subsection{Fonctionnelle utilisée}

Dans les deux expériences, l’énergie est de la forme
\[
J_{\mathrm{PCH}}(u;w,c)
=
\frac12\|u-w\|^2+\Phi_{\mathrm{PCH},c}(u),
\]
avec
\[
\Phi_{\mathrm{PCH},c}(u)
=
\lambda_{\mathrm{ctx}}\operatorname{dist}(u,\mathcal M_c)^2
+\lambda_{\mathrm{ridge}}\|u\|^2,
\qquad
\lambda_{\mathrm{ridge}}=\lambda_{\mathrm{ch}}+\lambda_{\mathrm{reg}}.
\]

Dans ex1, \(\lambda_{\mathrm{ridge}}=1.2\). Dans ex2, \(\lambda_{\mathrm{ridge}}\in\{0.2,0.8\}\).

\subsection{Définition explicite des routeurs comparés}

Pour éviter toute ambiguïté, les deux routeurs expérimentaux sont les suivants :

\[
\hat c_{\mathrm{energy}}(w)
=
\arg\min_{c\in\mathcal C}
J_{\mathrm{PCH}}(\Ppch(w;c);w,c),
\]
et
\[
\hat c_{\mathrm{proj}}(w)
=
\arg\min_{c\in\mathcal C}
\operatorname{dist}(w,\mathcal M_c)^2.
\]

Le premier est un diagnostic énergétique naturel mais fragile empiriquement ; le second est une heuristique géométrique plus robuste dans les runs testés. Le vrai contexte fourni par l’oracle sert de borne supérieure conditionnelle.

\subsection{Mesures géométriques}

On note
\[
d_c(u)=\operatorname{dist}(u,\mathcal M_c).
\]
Les réductions de distance rapportées dans les tableaux sont toujours des moyennes de \(d_c(w)\) avant purification et de \(d_c(\Ppch(w;c))\) après purification, avec le contexte oracle ou routé selon le cas.

\subsection{Résultats consolidés : ex1 (\texttt{20news20})}

Le régime ex1 est le plus exigeant sur le routage, car il comporte 20 contextes. Les résultats exacts sont :

\paragraph{Stabilité numérique et convergence.}
On observe :
\[
\texttt{energy\_monotonic\_violations}=0
\]
sur train et test, avec
\[
\texttt{mean\_fixed\_point\_residual}
=
4.086991364370892\times 10^{-17}\quad(\text{train}),
\qquad
4.2001606012501325\times 10^{-17}\quad(\text{test}).
\]
Les maxima atteignent respectivement
\[
9.999602279432701\times 10^{-17}\quad(\text{train}),
\qquad
1.3948925654884978\times 10^{-16}\quad(\text{test}).
\]

\paragraph{Loi locale de contraction.}
La norme de linéarisation locale vaut :
\[
\texttt{mean\_linearized\_norm}
=
0.22727272727485548\quad(\text{train}),
\qquad
0.2272727272751796\quad(\text{test}),
\]
avec
\[
\texttt{mean\_mu\_hat}
=
3.3999999999587978\quad(\text{train}),
\qquad
3.3999999999525228\quad(\text{test}),
\]
et des coefficients de variation extrêmement faibles :
\[
\texttt{cv\_mu\_hat}=1.6684207760963133\times 10^{-12}\quad(\text{train}),
\qquad
5.478711355470332\times 10^{-12}\quad(\text{test}).
\]

La loi observée est :
\[
\rho(\eta)=\frac{1}{1+\eta\mu_{\mathrm{eff}}}
\]
avec
\[
\mu_{\mathrm{eff}}=3.3999999999556603,
\qquad
\rho_{\mathrm{global}}=0.22727272727501754,
\qquad
\texttt{predicted\_rho\_from\_mu\_eff\_global}=0.22727272727501752,
\]
et erreur relative
\[
\texttt{relative\_rho\_error\_from\_mu\_eff\_global}
=
1.2212453270753656\times 10^{-16}.
\]

\paragraph{Ratios utiles.}
Les ratios utiles sont systématiquement en dessous de la linéarisation locale :
\[
\texttt{mean\_useful\_ratio}
=
0.21735222672793197\quad(\text{train}),
\qquad
0.2168894944937602\quad(\text{test}),
\]
avec des médianes
\[
0.21802581728549836\quad(\text{train}),
\qquad
0.21814003555311762\quad(\text{test}).
\]

\paragraph{Réduction géométrique.}
La distance moyenne au hull passe de :
\[
0.15324128763705894 \to 0.007461090086536279\quad(\text{train oracle}),
\qquad
0.16320475523763597 \to 0.007899542072163323\quad(\text{test oracle}).
\]
Sous routage par projection :
\[
0.13611497070592582 \to 0.006719888619855755\quad(\text{train}),
\qquad
0.13756447009654052 \to 0.006766477768671127\quad(\text{test}).
\]

La réduction de distance est donc
\[
\texttt{mean\_oracle\_distance\_reduction}
=
0.14578019755052263\quad(\text{train}),
\qquad
0.15530521316547266\quad(\text{test}),
\]
et
\[
\texttt{mean\_routed\_distance\_reduction}
=
0.12939508208607006\quad(\text{train}),
\qquad
0.1307979923278694\quad(\text{test}).
\]

\paragraph{Routage et tâche aval.}
L’accuracy du routage de contexte est :
\[
0.625\quad(\text{train}),\qquad 0.5859375\quad(\text{test}).
\]
Le probe brut vaut :
\[
\texttt{probe\_raw.accuracy}=0.5328125,
\qquad
\texttt{probe\_raw.macro\_f1}=0.5312542872902847.
\]
Le probe après PCH avec vrai contexte vaut :
\[
\texttt{probe\_oracle\_pch.accuracy}=0.971875,
\qquad
\texttt{probe\_oracle\_pch.macro\_f1}=0.9715533148487616.
\]
Le probe après PCH avec routage automatique vaut :
\[
\texttt{probe\_routed\_pch.accuracy}=0.5484375,
\qquad
\texttt{probe\_routed\_pch.macro\_f1}=0.5455746862795534.
\]
Les gains sont donc :
\[
\Delta\text{acc}_{\text{oracle}}=0.4390625,
\qquad
\Delta\text{acc}_{\text{routed}}=0.015625,
\]
\[
\Delta\text{macro-F1}_{\text{oracle}}=0.44029902755847694,
\qquad
\Delta\text{macro-F1}_{\text{routed}}=0.014320398989268734.
\]

\paragraph{Lecture scientifique de ex1.}
Le noyau PCH est numériquement très stable et fortement contractif, mais le routage à 20 contextes reste le goulot d’étranglement. Le gain conceptuel est maximal quand le vrai contexte est fourni ; il devient très faible quand le contexte est choisi automatiquement.

\subsection{Résultats consolidés : ex2 (\texttt{20news4}, 5 runs)}

Le sweep ex2 compare l’effet de \(\eta\), de \(\lambda_{\mathrm{ridge}}\) et du type de régularisation. Les cinq runs se regroupent en trois cas distincts, car deux paires sont strictement identiques numériquement :

\begin{itemize}
    \item \textbf{run0 = run1} exactement ;
    \item \textbf{run2 = run3} exactement ;
    \item \textbf{run4} est distinct.
\end{itemize}

Le tableau suivant synthétise les valeurs exactes utiles :

\begin{center}
\small
\begin{tabular}{lcccccc}
\hline
Run & \(\eta\) & \(\lambda_{\mathrm{ridge}}\) & régularisation & \(\mu_{\mathrm{eff}}\) & \(\rho\) & verdict local \\
\hline
0/1 & 1.0 & 0.2 & ridge / none & 1.3999999999918462 & 0.4166666666680822 & exact à \(10^{-16}\) \\
2/3 & 2.0 & 0.2 & ridge / none & 1.3999999999898165 & 0.2631578947382525 & exact à \(10^{-16}\) \\
4   & 1.0 & 0.8 & ridge & 2.5999999999824803 & 0.2777777777791296 & exact à \(10^{-16}\) \\
\hline
\end{tabular}
\end{center}

Les observations détaillées sont les suivantes.

\paragraph{Stabilité numérique.}
Dans tous les runs :
\[
\texttt{energy\_monotonic\_violations}=0,
\]
avec résidus de point fixe de l’ordre de \(10^{-17}\) à \(10^{-16}\). Par exemple :
\[
\texttt{mean\_fixed\_point\_residual}
=
1.7185635492005723\times 10^{-17}\quad(\text{run0/1 train}),
\]
\[
1.6670701135748873\times 10^{-17}\quad(\text{run0/1 test}),
\]
\[
2.4936548519777746\times 10^{-17}\quad(\text{run2/3 train}),
\]
\[
2.394770042412581\times 10^{-17}\quad(\text{run2/3 test}),
\]
\[
1.3401340310734358\times 10^{-17}\quad(\text{run4 train}),
\qquad
1.3931098026834653\times 10^{-17}\quad(\text{run4 test}).
\]

\paragraph{Loi de contraction.}
Les constantes observées sont :
\[
\rho=0.4166666666680822\quad(\eta=1,\ \lambda_{\mathrm{ridge}}=0.2),
\]
\[
\rho=0.2631578947382525\quad(\eta=2,\ \lambda_{\mathrm{ridge}}=0.2),
\]
\[
\rho=0.2777777777791296\quad(\eta=1,\ \lambda_{\mathrm{ridge}}=0.8).
\]
Elles coïncident avec
\[
\rho(\eta)=\frac{1}{1+\eta\mu_{\mathrm{eff}}}
\]
à la précision machine.

\paragraph{Paramètre effectif.}
On observe :
\[
\mu_{\mathrm{eff}}=1.3999999999918462\quad(\text{run0/1}),
\qquad
\mu_{\mathrm{eff}}=1.3999999999898165\quad(\text{run2/3}),
\qquad
\mu_{\mathrm{eff}}=2.5999999999824803\quad(\text{run4}).
\]
Autrement dit, \(\mu_{\mathrm{eff}}\) est stable à l’intérieur d’un réglage donné, mais pas universellement constant à travers les réglages.

\paragraph{Point clé non trivial.}
Les paires \texttt{ridge} / \texttt{none} sont strictement identiques dans les configurations 0/1 et 2/3. Cela signifie que, dans ces runs, le type de régularisation \emph{n’affecte pas} la dynamique observée. Le facteur réellement visible dans la constante locale est \(\lambda_{\mathrm{ridge}}\), pas l’étiquette \texttt{ridge} elle-même.

\paragraph{Hypothèse empirique forte suggérée par les chiffres.}
Les trois runs distincts sont compatibles avec :
\[
\mu_{\mathrm{eff}}
\approx
1+2\lambda_{\mathrm{ridge}}.
\]



En effet :
\[
1+2\times 0.2 = 1.4,
\qquad
1+2\times 0.8 = 2.6.
\]
Par conséquent,
\[
\rho(\eta)
\approx
\frac{1}{1+\eta(1+2\lambda_{\mathrm{ridge}})}.
\]
Cette loi est \emph{empiriquement très forte} dans le périmètre testé, mais elle n’est pas encore un théorème universel.

Cependant, le nombre de valeurs testées de lambda_{\mathrm{ridge}} reste insuffisant pour établir cette relation comme loi générale.

\paragraph{Routage et tâche aval.}
Pour run0/1 et run2/3 :
\[
\texttt{probe\_raw.accuracy}=0.75,\qquad \texttt{probe\_raw.macro\_f1}=0.751195474136017,
\]
\[
\texttt{probe\_oracle\_pch.accuracy}=0.8515625,\qquad \texttt{probe\_oracle\_pch.macro\_f1}=0.851615787917469,
\]
\[
\texttt{probe\_routed\_pch.accuracy}=0.796875,\qquad \texttt{probe\_routed\_pch.macro\_f1}=0.7977103235961343.
\]
Le routage de contexte vaut :
\[
0.84375\ (\text{train}),\ 0.8359375\ (\text{test}).
\]

Pour run4 :
\[
\texttt{probe\_oracle\_pch.accuracy}=0.84375,\qquad \texttt{probe\_oracle\_pch.macro\_f1}=0.8434727880769546,
\]
\[
\texttt{probe\_routed\_pch.accuracy}=0.796875,\qquad \texttt{probe\_routed\_pch.macro\_f1}=0.797720286885246.
\]

\paragraph{Réduction géométrique.}
Pour run0/1 :
\[
0.1346794788494753 \to 0.038787138578643535\quad(\text{oracle train}),
\]
\[
0.13667472169476488 \to 0.03918071848335118\quad(\text{oracle test}),
\]
et sous routage :
\[
0.11908961214321748 \to 0.03347963381159686\quad(\text{train}),
\]
\[
0.12194849604687594 \to 0.034048417681437154\quad(\text{test}).
\]

Pour run4 :
\[
0.1346794788494753 \to 0.021216944702244885\quad(\text{oracle train}),
\]
\[
0.13667472169476488 \to 0.021380746092661156\quad(\text{oracle test}),
\]
et sous routage :
\[
0.11908961214321748 \to 0.018888601997152858\quad(\text{train}),
\]
\[
0.12194849604687594 \to 0.01889335597174899\quad(\text{test}).
\]

\paragraph{Verdict du sweep.}
Le verdict automatique donne :
\[
\texttt{supported}=false,
\qquad
\texttt{mu\_eff\_cv}=0.32722946012180587.
\]
Ce refus ne contredit pas la loi locale observée à l’intérieur de chaque run ; il réfute seulement l’hypothèse plus forte d’une \emph{constante universelle unique} \(\mu_{\mathrm{eff}}\) commune à tous les hyperparamètres du sweep.

\subsection{Interprétation rigoureuse}

Les données établissent de manière robuste les points suivants :

\begin{enumerate}
    \item \textbf{Cas 1 est confirmé.}
    Le noyau PCH est bien posé, le minimiseur existe numériquement, l’itération proximale converge, l’énergie décroît sans violation observée, et les résidus de point fixe restent à l’échelle \(10^{-17}\) à \(10^{-16}\).

    \item \textbf{La loi locale de contraction est exacte à la précision machine.}
    Dans ex1 :
    \[
    \rho_{\mathrm{lin}} \approx 0.22727272727501754,
    \qquad
    \mu_{\mathrm{eff}} \approx 3.3999999999556603.
    \]
    Dans ex2 :
    \[
    \rho = \frac{1}{1+\eta\mu_{\mathrm{eff}}}
    \]
    est vérifiée avec des erreurs relatives de l’ordre de \(10^{-16}\).

    \item \textbf{Cas 2 est fortement supporté.}
    La réduction de distance au hull admissible est systématique et forte.
    En ex1, l’oracle fait passer la distance moyenne d’environ \(0.1532\) à \(0.00746\) en train, et de \(0.1632\) à \(0.00790\) en test.
    En ex2, les réductions restent nettes, mais plus modestes, ce qui montre que la géométrie du contexte, la qualité du routage et la granularité du corpus dominent le gain aval.

    \item \textbf{Le routage est le vrai goulot.}
    En ex1, le routage automatique atteint seulement
    \[
    0.625\ (\text{train}),\quad 0.5859375\ (\text{test}),
    \]
    alors que l’oracle donne
    \[
    0.971875
    \]
    sur le probe.
    En ex2, avec seulement 4 contextes, le routage devient beaucoup moins pénalisant :
    \[
    0.84375\ (\text{train}),\quad 0.8359375\ (\text{test}),
    \]
    ce qui confirme que la difficulté croît principalement avec le nombre de contextes.

    \item \textbf{Le sweep réfute l’universalité stricte de \(\mu_{\mathrm{eff}}\), mais pas la loi locale.}
    Les valeurs \(\mu_{\mathrm{eff}}=1.4\) et \(\mu_{\mathrm{eff}}=2.6\) montrent que \(\mu_{\mathrm{eff}}\) dépend des hyperparamètres, en particulier de \(\lambda_{\mathrm{ridge}}\).
    En revanche, à l’intérieur de chaque run, la relation résolvante est exacte à la précision machine.
\end{enumerate}

\subsection{Clarification sur le statut du choix de contexte}

La difficulté observée sur le choix de contexte ne signifie pas que la théorie est fausse. Elle signifie que la théorie conditionnelle du noyau PCH est validée, tandis que la règle de sélection automatique du contexte n’est pas encore certifiée comme loi universelle. Autrement dit, le bloc PCH fonctionne ; le routeur universel n’est pas encore démontré.

\subsection{Points clés supplémentaires qui méritent d’être conservés}

Trois points ressortent clairement des chiffres et peuvent passer inaperçus dans une lecture trop rapide :

\begin{itemize}
    \item \textbf{L’effet de \texttt{ridge} versus \texttt{none} est nul dans les runs appariés.}
    Les runs 0/1 et 2/3 sont identiques numériquement.
    Cela signifie que, dans ce protocole, le simple changement d’étiquette de régularisation ne modifie pas la dynamique observée.

    \item \textbf{La variation de \(\eta\) agit exactement comme une loi de résolvant.}
    Quand \(\mu_{\mathrm{eff}}=1.4\), on passe de
    \[
    0.4166666666680822 \quad \text{à} \quad 0.2631578947382525
    \]
    lorsque \(\eta\) passe de \(1\) à \(2\), sans changer \(\mu_{\mathrm{eff}}\).

    \item \textbf{La variation de \(\lambda_{\mathrm{reg}}\) déplace \(\mu_{\mathrm{eff}}\) de façon affine.}
    On passe de \(1.4\) à \(2.6\) quand \(\lambda_{\mathrm{reg}}\) passe de \(0.2\) à \(0.8\), ce qui est compatible avec
    \[
    \mu_{\mathrm{eff}} \approx 1 + 2\lambda_{\mathrm{ridge}}.
    \]
\end{itemize}

\subsection{Variabilité, robustesse et limites}

Les écarts-types reportés dans ex2 sont minuscules à l’intérieur de chaque réglage :
\[
\texttt{std\_mu\_hat}=2.7706473529209935\times 10^{-12}\ \text{à}\ 6.2571818240720185\times 10^{-12}\quad(\text{train}),
\]
\[
\texttt{std\_mu\_hat}=2.818001749473456\times 10^{-12}\ \text{à}\ 4.122917216587002\times 10^{-12}\quad(\text{test}),
\]
ce qui confirme la stabilité locale du phénomène observé dans chaque run. En revanche, l’agrégation sur plusieurs réglages fait apparaître une variation réelle de \(\mu_{\mathrm{eff}}\), donc une dépendance aux hyperparamètres.

Cette campagne demeure limitée à un petit nombre de familles expérimentales :
un encodeur fixé par campagne, une réduction PCA fixée, deux corpus principaux (\texttt{20news20} et \texttt{20news4}), des contextes convexes définis par calibration, et des routeurs heuristiques. La généralisation à d’autres encodeurs, à d’autres corpus, à d’autres régularisations et à d’autres modes de routage reste un problème ouvert.


\subsection{Reproductibilité}

Toutes les statistiques reportées proviennent des sorties
générées automatiquement par le pipeline expérimental.
Les valeurs numériques sont reproduites telles qu'elles
apparaissent dans les logs d'exécution et ne résultent
d'aucun arrondi manuel, sauf indication contraire.

\subsection{Conséquences conceptuelles, théoriques, opérationnelles et pratiques}

\paragraph{Conséquence conceptuelle.}
Le noyau PCH ne se comporte pas comme une heuristique libre ; il se comporte comme un opérateur proximal géométrique à structure locale très rigide.
La loi locale observée n’est pas une simple corrélation : elle est numériquement exacte à la précision machine sur les runs reportés.

\paragraph{Conséquence théorique.}
La hiérarchie correcte est désormais la suivante :
\[
\text{Cas 1 (convergence proximale)} \;\Rightarrow\; \text{confirmé},
\]
\[
\text{Cas 2 (réduction géométrique locale)} \;\Rightarrow\; \text{fortement supporté},
\]
\[
\text{Cas 3 local (forme résolvante avec }H_{\mathrm{eff}}\text{)} \;\Rightarrow\; \text{très fortement supporté},
\]
mais
\[
\text{Cas 3 universel (même }H_{\mathrm{eff}}\text{ pour tous les paramètres)} \;\Rightarrow\; \text{non démontré}.
\]

Autrement dit, le bon énoncé n’est pas :
\[
D\Tpch(u^\star)=(I+\eta H_{\mathrm{eff}})^{-1}
\quad\text{avec un }H_{\mathrm{eff}}\text{ universel},
\]
mais plutôt :
\[
D\Tpch(u^\star)\approx (I+\eta H_{\mathrm{eff}}(\lambda_{\mathrm{ridge}},\mathcal M_c,E_\theta,d))^{-1},
\]
avec une loi locale extrêmement stable, mais dépendante du régime expérimental.

\paragraph{Conséquence opérationnelle.}
Le routage est le principal facteur limitant dans les régimes à grand nombre de contextes.
À 20 contextes, le noyau reste stable, mais la sélection du bon contexte s’effondre relativement par rapport à l’oracle.
À 4 contextes, le routage devient nettement plus fiable.

\paragraph{Conséquence pratique.}
Le protocole gagnant n’est pas d’augmenter sans fin \(\texttt{max\_iters}\), mais de stabiliser le routage et la géométrie des contexts.
Les mesures montrent que :
\[
\text{qualité du routeur} \gg \text{nombre d’itérations}
\]
une fois la convergence atteinte.

\paragraph{Conséquence pour la lecture du modèle.}
PCH ne doit pas être présenté comme un classifieur universel.
Le bloc PCH agit d’abord comme une \emph{purification contextuelle exacte} :
il rapproche l’embedding d’un ensemble admissible, améliore parfois le probe, mais n’engendre pas à lui seul la séparabilité.

\subsection{Verdict scientifique}

Le verdict exact et intégral est le suivant :

\begin{itemize}
    \item \(\boxed{\text{Cas 1 est confirmé}}\) : convergence, stabilité, monotonie énergétique, point fixe.
    \item \(\boxed{\text{Cas 2 est fortement supporté}}\) : réduction géométrique au hull admissible.
    \item \(\boxed{\text{La loi locale de contraction est confirmée à la précision machine}}\) :
    \[
    \rho(\eta)=\frac{1}{1+\eta\mu_{\mathrm{eff}}}.
    \]
    \item \(\boxed{\text{L’hypothèse d’un } \mu_{\mathrm{eff}} \text{ universel unique est réfutée par le sweep}}\).
    \item \(\boxed{\text{La meilleure lecture actuelle est une loi de résolvant locale dépendant des hyperparamètres}}\).
\end{itemize}

La découverte la plus importante n’est donc pas seulement que PCH converge.
C’est que, dans les régimes testés, la dynamique locale semble obéir à une loi de type résolvant,
avec
\[
\mu_{\mathrm{eff}} \approx 1 + 2\lambda_{\mathrm{ridge}},
\]
vérifiée à la précision machine sur les runs reportés, tandis que le routage reste la partie fragile du système.




%%%%%%%%%%%%%%%%%%%%%%%%%%%%%%%%%%%%%%%%%%%%%%%%%%%%%%%%%%%%%%%%%%%%%%%%%%%%%%%%%%%%%%%%%%%%%%%%%%%%%%%%%%%%%%%%%%%%%%%%%%%%%%%%%%%%%%%%%%%%%%%%%%%%%%%%%%%%%%%%%%%%%%%%%%%%%%%%%%%%%%%%%%%%%%%%%%%%%%%%%%%%%%%%%%%%%%%%%%%%%%%%%%%%%%%%%%%%%%%%%%%%%%%%%%%%%%%%%%%%%%%%%%%%%%%%%%%%%%%%%%%demonstration cas 3 local :

\section{Analyse locale du projection hull et loi résolvante du noyau PCH}

\subsection{Cadre et hypothèses locales}

Soit
\[
M=\operatorname{conv}(x_1,\dots,x_m)\subset\mathbb R^d
\]
un polytope convexe compact non vide. On note
\[
\Pi_M:\mathbb R^d\to M
\]
la projection euclidienne sur $M$, et
\[
d_M(u)=\operatorname{dist}(u,M)=\|u-\Pi_M(u)\|
\]
la distance euclidienne à $M$.

On considère un point $u^\star\in\mathbb R^d$ tel que
\[
p^\star:=\Pi_M(u^\star)
\]
appartienne au relatif intérieur d'une face $F$ de $M$.

On suppose en outre que la face active reste stable dans un voisinage de $u^\star$, c'est-à-dire qu'il existe un voisinage $U$ de $u^\star$ tel que, pour tout $u\in U$, la projection $\Pi_M(u)$ reste dans le relatif intérieur de la même face $F$.

On note
\[
T_F:=\operatorname{lin}(F-F)
\qquad\text{et}\qquad
N_F:=T_F^\perp,
\]
et on désigne par
\[
P_T:\mathbb R^d\to T_F,\qquad P_N:\mathbb R^d\to N_F
\]
les projecteurs orthogonaux correspondants.

\subsection{Lemme 1 : forme locale de la projection convexe}

\begin{lemma}
Sous les hypothèses ci-dessus, la projection $\Pi_M$ est différentiable en $u^\star$ et
\[
D\Pi_M(u^\star)=P_T.
\]
\end{lemma}

\begin{proof}
Puisque $p^\star$ appartient au relatif intérieur de la face $F$ et que la face active reste inchangée dans un voisinage de $u^\star$, la projection sur $M$ coïncide localement avec la projection orthogonale sur l'espace affine $\operatorname{aff}(F)$, suivie de la restriction au voisinage où le point projeté reste dans $F$.

Or la projection orthogonale sur un sous-espace affine est une application affine, donc différentiable partout, dont la différentielle est le projecteur orthogonal sur la direction linéaire de ce sous-espace. Cette direction est précisément $T_F$.

Ainsi,
\[
D\Pi_M(u^\star)=P_T.
\]
\end{proof}

\subsection{Lemme 2 : Hessien local de la distance carrée}

On considère la fonction
\[
\varphi(u)=d_M(u)^2=\|u-\Pi_M(u)\|^2.
\]

\begin{lemma}
Si $\Pi_M$ est différentiable en $u$, alors
\[
\nabla^2 \varphi(u)=2\bigl(I-D\Pi_M(u)\bigr).
\]
En particulier, sous les hypothèses locales précédentes,
\[
\nabla^2 \varphi(u^\star)=2(I-P_T)=2P_N.
\]
\end{lemma}

\begin{proof}
Le fait classique de l'analyse convexe est que
\[
\nabla\!\left(\frac12 d_M(u)^2\right)=u-\Pi_M(u).
\]
En multipliant par $2$, on obtient
\[
\nabla \varphi(u)=2(u-\Pi_M(u)).
\]
Si $\Pi_M$ est différentiable en $u$, on dérive cette identité :
\[
\nabla^2\varphi(u)=2I-2D\Pi_M(u)=2\bigl(I-D\Pi_M(u)\bigr).
\]
Au point $u^\star$, le lemme précédent donne $D\Pi_M(u^\star)=P_T$, donc
\[
\nabla^2\varphi(u^\star)=2(I-P_T)=2P_N.
\]
\end{proof}

\subsection{Structure spectrale de la projection locale}

\begin{proposition}
Sous les hypothèses locales précédentes, le spectre de $D\Pi_M(u^\star)$ est exactement
\[
\operatorname{Spec}(D\Pi_M(u^\star))
=
\{1\}^{\dim(T_F)}\cup \{0\}^{d-\dim(T_F)}.
\]
Par conséquent,
\[
\operatorname{Spec}(\nabla^2\varphi(u^\star))
=
\{0\}^{\dim(T_F)}\cup \{2\}^{d-\dim(T_F)}.
\]
\end{proposition}

\begin{proof}
Le lemme 1 montre que $D\Pi_M(u^\star)=P_T$. Or $P_T$ est un projecteur orthogonal, donc
\[
P_T^2=P_T,\qquad P_T^\top=P_T.
\]
Ses valeurs propres vérifient $\lambda^2=\lambda$, donc
\[
\lambda\in\{0,1\}.
\]
La multiplicité de la valeur propre $1$ est $\dim(T_F)$, et celle de $0$ est $d-\dim(T_F)$.

Comme
\[
\nabla^2\varphi(u^\star)=2(I-P_T),
\]
les valeurs propres correspondantes sont $0$ sur $T_F$ et $2$ sur $N_F$.
\end{proof}

\subsection{Fonctionnelle PCH locale}

On considère l'énergie PCH
\[
J(u;w,c)
=
\frac12\|u-w\|^2
+
\lambda_{\mathrm{ctx}}\,d_M(u)^2
+
\lambda_{\mathrm{ridge}}\|u\|^2.
\]

Sous les hypothèses locales ci-dessus, $J(\cdot;w,c)$ est de classe $C^2$ près de $u^\star$ et
\[
\nabla^2 J(u^\star;w,c)
=
I+2\lambda_{\mathrm{ridge}}I+2\lambda_{\mathrm{ctx}}(I-P_T).
\]
Autrement dit,
\[
\nabla^2 J(u^\star;w,c)
=
(1+2\lambda_{\mathrm{ridge}})I+2\lambda_{\mathrm{ctx}}(I-P_T).
\]

\begin{proof}
On dérive terme à terme :
\[
\nabla^2\Big(\frac12\|u-w\|^2\Big)=I,\qquad
\nabla^2\big(\lambda_{\mathrm{ridge}}\|u\|^2\big)=2\lambda_{\mathrm{ridge}}I,
\]
et, par le lemme précédent,
\[
\nabla^2\big(\lambda_{\mathrm{ctx}}d_M(u)^2\big)
=
2\lambda_{\mathrm{ctx}}(I-P_T).
\]
La somme donne la formule annoncée.
\end{proof}

\subsection{Jacobian local de l'opérateur proximal}

On définit l'opérateur proximal exact
\[
T_\eta(z;w,c)
=
\arg\min_u
\left\{
\frac12\|u-w\|^2
+
\lambda_{\mathrm{ctx}}\,d_M(u)^2
+
\lambda_{\mathrm{ridge}}\|u\|^2
+
\frac1{2\eta}\|u-z\|^2
\right\}.
\]

Soit $u^\star$ un point fixe local, c'est-à-dire
\[
u^\star=T_\eta(u^\star;w,c).
\]

\begin{theorem}[Loi résolvante locale exacte]
Sous les hypothèses précédentes,
\[
DT_\eta(u^\star;w,c)
=
\bigl(I+\eta \nabla^2J(u^\star;w,c)\bigr)^{-1}.
\]
En particulier, le spectre de $DT_\eta(u^\star;w,c)$ est donné par
\[
\operatorname{Spec}(DT_\eta(u^\star;w,c))
=
\left\{
\frac1{1+\eta(1+2\lambda_{\mathrm{ridge}})}
\right\}^{\dim(T_F)}
\cup
\left\{
\frac1{1+\eta(1+2\lambda_{\mathrm{ridge}}+2\lambda_{\mathrm{ctx}})}
\right\}^{d-\dim(T_F)}.
\]
\end{theorem}

\begin{proof}
Posons
\[
\Phi(u)=\frac12\|u-w\|^2+\lambda_{\mathrm{ctx}}\,d_M(u)^2+\lambda_{\mathrm{ridge}}\|u\|^2.
\]
L'opérateur proximal s'écrit
\[
T_\eta(z)=\arg\min_u\left\{\Phi(u)+\frac1{2\eta}\|u-z\|^2\right\}.
\]
La condition d'optimalité s'écrit
\[
\nabla\Phi(u)+\frac1\eta(u-z)=0.
\]
Au point fixe local $(u^\star,z^\star=u^\star)$, on applique le théorème des fonctions implicites. En différentiant par rapport à $z$, on obtient
\[
\left(\nabla^2\Phi(u^\star)+\frac1\eta I\right)DT_\eta(u^\star)
=
\frac1\eta I.
\]
Donc
\[
DT_\eta(u^\star)
=
\left(I+\eta\nabla^2\Phi(u^\star)\right)^{-1}.
\]
Or $\nabla^2\Phi(u^\star)=\nabla^2J(u^\star;w,c)$, ce qui donne la formule annoncée.

Enfin, comme
\[
\nabla^2J(u^\star;w,c)
=
(1+2\lambda_{\mathrm{ridge}})I+2\lambda_{\mathrm{ctx}}(I-P_T),
\]
les directions tangentes à la face active ont valeur propre
\[
\mu_T=1+2\lambda_{\mathrm{ridge}},
\]
et les directions normales ont valeur propre
\[
\mu_N=1+2\lambda_{\mathrm{ridge}}+2\lambda_{\mathrm{ctx}}.
\]
On obtient donc les valeurs propres de $DT_\eta(u^\star;w,c)$ en inversant $1+\eta\mu_T$ et $1+\eta\mu_N$.
\end{proof}

\subsection{Conséquence : constante de contraction observée}

Comme
\[
\mu_N\ge \mu_T,
\]
le rayon spectral est
\[
\rho\bigl(DT_\eta(u^\star;w,c)\bigr)
=
\frac1{1+\eta(1+2\lambda_{\mathrm{ridge}})}.
\]

Autrement dit, la constante de contraction locale observée ne dépend pas de la géométrie fine de la face active, mais seulement du terme quadratique injecté par
\[
\frac12\|u-w\|^2+\lambda_{\mathrm{ridge}}\|u\|^2,
\]
la contribution contextuelle n'affectant que les directions normales.

Ceci explique exactement les valeurs expérimentales :
\[
\lambda_{\mathrm{ridge}}=0.2 \quad\Longrightarrow\quad
\rho=\frac1{1+1.4}=\frac5{12}\approx 0.4166666667,
\]
\[
\lambda_{\mathrm{ridge}}=0.2,\ \eta=2 \quad\Longrightarrow\quad
\rho=\frac1{1+2\times 1.4}=\frac5{19}\approx 0.2631578947,
\]
\[
\lambda_{\mathrm{ridge}}=0.8 \quad\Longrightarrow\quad
\rho=\frac1{1+2.6}=\frac5{18}\approx 0.2777777778.
\]

\subsection{Ce qui est démontré et ce qui ne l'est pas}

Les résultats ci-dessus démontrent rigoureusement :

\begin{enumerate}
\item la formule locale exacte
\[
\nabla^2 d_M(u)^2=2(I-D\Pi_M(u));
\]

\item la structure spectrale de la projection convexe au voisinage d'une face active stable :
\[
\operatorname{Spec}(D\Pi_M(u^\star))\subset\{0,1\};
\]

\item la forme résolvante locale exacte du proximal :
\[
DT_\eta(u^\star;w,c)
=
\bigl(I+\eta\nabla^2J(u^\star;w,c)\bigr)^{-1};
\]

\item la loi locale de contraction :
\[
\rho
=
\frac1{1+\eta(1+2\lambda_{\mathrm{ridge}})}.
\]
\end{enumerate}

En revanche, la propriété suivante n'est pas démontrée par l'analyse locale seule :

\[
\text{l'universalité du même }H_{\mathrm{eff}}\text{ à travers tous les corpus, encodeurs et PCA.}
\]

Cette dernière affirmation requerrait un théorème supplémentaire de concentration géométrique des faces actives, ou une loi asymptotique de stabilité du projecteur tangent à travers la famille des hulls réels.

\subsection{Version Cas 3 local}

Sous les hypothèses locales de régularité de face active, on peut donc énoncer le Cas 3 local de manière rigoureuse :

\[
\boxed{
DT_\eta(u^\star;w,c)
=
\bigl(I+\eta H_{\mathrm{eff}}(c,w)\bigr)^{-1}
}
\]
avec
\[
H_{\mathrm{eff}}(c,w)
=
(1+2\lambda_{\mathrm{ridge}})I+2\lambda_{\mathrm{ctx}}(I-P_T).
\]

Dans ce sens précis, le noyau PCH possède bien une loi de type résolvant, et la constante de contraction observée n'est pas accidentelle : elle est la conséquence exacte de la structure convexe locale du problème.

\subsection{Remarque finale}

Si la face active est un singleton, alors $P_T=0$ et la formule devient
\[
DT_\eta(u^\star;w,c)
=
\frac1{1+\eta(1+2\lambda_{\mathrm{ridge}}+2\lambda_{\mathrm{ctx}})}I.
\]
Ce cas est particulier et ne décrit pas les runs expérimentaux consolidés, dans lesquels la direction tangentielle existe et gouverne le rayon spectral local.








%%%%%%%%%%%%%%%%%%%%%%%%%%%%%%%%%%%%%%%%%%%%%%%%%%%%%%%%%%%%%%%%%%%%%%%%%%%%%%%%%%%%%%%%%%%%%%%%%%%%%%%%%%%%%%%%%%%%%%%%%%%%%%%%%%%%%%%%%%%%%%%%%%%%%%%%%%%%%%%%%%%%%%%%%%%%%%%%%%%%%%%%%%%%%%%%%%%%%%%%%%%%%%%%%%%%%%%%%%%%%%%%%%%%%%%%%%%%%%%%%%%%%%%%%%%%%%%%%%%%%%%%%%%%%%%%%%%%%%%%%%%demonstration cas 3 universel :

\section{Cas 3 universel sous hypothèse de concentration géométrique}

\subsection{Cadre}

Soit une famille de problèmes PCH indexée par \(\alpha\), où \(\alpha\) encode le corpus, l'encodeur, la PCA, la dimension latente et les hyperparamètres. Pour chaque \(\alpha\), on considère un contexte \(c\), un ensemble admissible convexe
\[
M_\alpha \subset \mathbb R^{d_\alpha},
\]
et l'énergie locale
\[
J_\alpha(u;w,c)
=
\frac12\|u-w\|^2
+
\lambda_{\mathrm{ctx},\alpha}\,\dist(u,M_\alpha)^2
+
\lambda_{\mathrm{ridge},\alpha}\,\|u\|^2.
\]
On note
\[
T_{\eta,\alpha}(z;w,c)
=
\arg\min_u
\left\{
\frac12\|u-w\|^2
+
\lambda_{\mathrm{ctx},\alpha}\,\dist(u,M_\alpha)^2
+
\lambda_{\mathrm{ridge},\alpha}\,\|u\|^2
+
\frac{1}{2\eta}\|u-z\|^2
\right\}.
\]

\subsection{Hypothèse de concentration géométrique universelle}

On suppose qu'il existe un projecteur orthogonal déterministe \(P_0\) et une suite \(\varepsilon_\alpha\to 0\) tels que, pour tout \(\alpha\), au point fixe local pertinent \(u_\alpha^\star\),
\[
\|D\Pi_{M_\alpha}(u_\alpha^\star)-P_0\|_{\mathrm{op}}
\le \varepsilon_\alpha,
\]
et que cette convergence est uniforme sur la famille considérée.

On suppose en outre que les coefficients convergent vers des constantes universelles :
\[
\lambda_{\mathrm{ctx},\alpha}\to \lambda_{\mathrm{ctx},\star},
\qquad
\lambda_{\mathrm{ridge},\alpha}\to \lambda_{\mathrm{ridge},\star}.
\]

\subsection{Lemme 1 : Hessien de la distance carrée}

\begin{lemma}
Si \(\Pi_{M_\alpha}\) est différentiable en \(u\), alors
\[
\nabla^2 \dist(u,M_\alpha)^2
=
2\bigl(I-D\Pi_{M_\alpha}(u)\bigr).
\]
\end{lemma}

\begin{proof}
On utilise l'identité convexe classique
\[
\nabla\!\left(\frac12\dist(u,M_\alpha)^2\right)=u-\Pi_{M_\alpha}(u).
\]
En dérivant,
\[
\nabla^2\!\left(\frac12\dist(u,M_\alpha)^2\right)
=
I-D\Pi_{M_\alpha}(u),
\]
d'où
\[
\nabla^2\dist(u,M_\alpha)^2
=
2\bigl(I-D\Pi_{M_\alpha}(u)\bigr).
\]
\end{proof}

\subsection{Lemme 2 : concentration du Hessien géométrique}

\begin{lemma}
Sous l'hypothèse de concentration géométrique,
\[
\nabla^2 \dist(u_\alpha^\star,M_\alpha)^2
\longrightarrow
2(I-P_0)
\]
en norme opératorielle.
\end{lemma}

\begin{proof}
Par le lemme précédent,
\[
\nabla^2 \dist(u_\alpha^\star,M_\alpha)^2
=
2\bigl(I-D\Pi_{M_\alpha}(u_\alpha^\star)\bigr).
\]
Donc
\[
\left\|
\nabla^2 \dist(u_\alpha^\star,M_\alpha)^2
-
2(I-P_0)
\right\|_{\mathrm{op}}
\le
2\|D\Pi_{M_\alpha}(u_\alpha^\star)-P_0\|_{\mathrm{op}}
\le
2\varepsilon_\alpha\to 0.
\]
\end{proof}

\subsection{Théorème 1 : Hessien effectif universel}

\begin{theorem}
Sous l'hypothèse de concentration géométrique universelle,
\[
\nabla^2 J_\alpha(u_\alpha^\star;w,c)
\longrightarrow
H_\star
\]
avec
\[
H_\star
=
(1+2\lambda_{\mathrm{ridge},\star})I
+
2\lambda_{\mathrm{ctx},\star}(I-P_0).
\]
\end{theorem}

\begin{proof}
On a
\[
\nabla^2 J_\alpha(u_\alpha^\star;w,c)
=
I
+
2\lambda_{\mathrm{ridge},\alpha}I
+
\lambda_{\mathrm{ctx},\alpha}\,\nabla^2 \dist(u_\alpha^\star,M_\alpha)^2.
\]
En appliquant le lemme précédent,
\[
\nabla^2 J_\alpha(u_\alpha^\star;w,c)
=
I
+
2\lambda_{\mathrm{ridge},\alpha}I
+
2\lambda_{\mathrm{ctx},\alpha}\bigl(I-D\Pi_{M_\alpha}(u_\alpha^\star)\bigr).
\]
Par convergence de \(\lambda_{\mathrm{ridge},\alpha}\), \(\lambda_{\mathrm{ctx},\alpha}\), et \(D\Pi_{M_\alpha}(u_\alpha^\star)\), on obtient
\[
\nabla^2 J_\alpha(u_\alpha^\star;w,c)
\to
(1+2\lambda_{\mathrm{ridge},\star})I
+
2\lambda_{\mathrm{ctx},\star}(I-P_0)
=
H_\star.
\]
\end{proof}

\subsection{Théorème 2 : forme résolvante du Jacobien du prox}

\begin{theorem}
Sous les hypothèses précédentes,
\[
DT_{\eta,\alpha}(u_\alpha^\star;w,c)
\longrightarrow
(I+\eta H_\star)^{-1}
\]
en norme opératorielle.
\end{theorem}

\begin{proof}
Le point fixe \(u_\alpha^\star\) satisfait la condition d'optimalité du proximal :
\[
0
=
\nabla J_\alpha(u_\alpha^\star;w,c)
+
\frac1\eta(u_\alpha^\star-z).
\]
Par le théorème des fonctions implicites, la différentielle de la solution par rapport à \(z\) est
\[
DT_{\eta,\alpha}(u_\alpha^\star;w,c)
=
\bigl(I+\eta \nabla^2 J_\alpha(u_\alpha^\star;w,c)\bigr)^{-1}.
\]
Comme \(\nabla^2 J_\alpha(u_\alpha^\star;w,c)\to H_\star\), on obtient
\[
DT_{\eta,\alpha}(u_\alpha^\star;w,c)
\to
(I+\eta H_\star)^{-1}.
\]
\end{proof}

\subsection{Conséquence spectrale}

Comme \(P_0\) est un projecteur orthogonal, son spectre est contenu dans \(\{0,1\}\). Si \(r=\mathrm{rank}(P_0)\), alors
\[
\operatorname{Spec}(H_\star)
=
\left\{1+2\lambda_{\mathrm{ridge},\star}\right\}^{r}
\cup
\left\{1+2\lambda_{\mathrm{ridge},\star}+2\lambda_{\mathrm{ctx},\star}\right\}^{d-r}.
\]
Par conséquent,
\[
\operatorname{Spec}\bigl((I+\eta H_\star)^{-1}\bigr)
=
\left\{
\frac{1}{1+\eta(1+2\lambda_{\mathrm{ridge},\star})}
\right\}^{r}
\cup
\left\{
\frac{1}{1+\eta(1+2\lambda_{\mathrm{ridge},\star}+2\lambda_{\mathrm{ctx},\star})}
\right\}^{d-r}.
\]
Le rayon spectral local est donc
\[
\rho_\star
=
\frac{1}{1+\eta(1+2\lambda_{\mathrm{ridge},\star})}.
\]

\subsection{Cas 3 universel}

\begin{theorem}[Cas 3 universel]
Sous l'hypothèse de concentration géométrique universelle des faces actives, la dynamique locale du noyau PCH est asymptotiquement gouvernée par un opérateur résolvant universel :
\[
DT_{\eta,\alpha}(u_\alpha^\star;w,c)
\to
(I+\eta H_\star)^{-1},
\]
où
\[
H_\star
=
(1+2\lambda_{\mathrm{ridge},\star})I
+
2\lambda_{\mathrm{ctx},\star}(I-P_0),
\]
et où le rayon spectral local est universellement donné par
\[
\rho_\star
=
\frac{1}{1+\eta(1+2\lambda_{\mathrm{ridge},\star})}.
\]
\end{theorem}

\begin{proof}
C'est la réunion du Théorème 1 et du Théorème 2.
\end{proof}

\subsection{Lecture finale}

Le Cas 3 universel est donc un théorème \emph{conditionnel} :
la seule hypothèse non triviale est la concentration géométrique uniforme du différentiel de projection autour d'un projecteur limite \(P_0\).  
Une fois cette hypothèse admise, la forme résolvante du Jacobien du prox suit mécaniquement, et la loi spectrale universelle devient une conséquence directe.




%%%%%%%%%%%%%%%%%%%%%%%%%%%%%%%%%%%%%%%%%%%%%%%%%%%%%%%%%%%%%%%%%%%%%%%%%%%%%%%%%%%%%%%%%%%%%%%%%%%%%%%%%%%%%%%%%%%%%%%%%%%%%%%%%%%%%%%%%%%%%%%%%%%%%%%%%%%%%%%%%%%%%%%%%%%%%%%%%%%%%%%%%%%%%%%%%%%%%%%%%%%%%%%%%%%%%%%%%%%%%%%%%%%%%%%%%%%%%%%%%%%%%%%%%%%%%%%%%%%%%%%%%%%%%%%%%%%%%%%%%%% cadre proba

\section{Cas 3 universel sous hypothèse de concentration géométrique}

\subsection{Cadre probabiliste et objet limite}

Soit \(\alpha\) un index de modèle, qui encode le corpus, l'encodeur, la PCA, la dimension latente et le choix des hyperparamètres. Pour chaque \(\alpha\), on considère un contexte \(c\), un ensemble admissible convexe
\[
M_\alpha \subset \mathbb R^{d_\alpha},
\]
et l'énergie locale
\[
J_\alpha(u;w,c)
=
\frac12\|u-w\|^2
+
\lambda_{\mathrm{ctx},\alpha}\,\dist(u,M_\alpha)^2
+
\lambda_{\mathrm{ridge},\alpha}\,\|u\|^2.
\]
On définit le proximal associé par
\[
T_{\eta,\alpha}(z;w,c)
=
\arg\min_u
\left\{
\frac12\|u-w\|^2
+
\lambda_{\mathrm{ctx},\alpha}\,\dist(u,M_\alpha)^2
+
\lambda_{\mathrm{ridge},\alpha}\,\|u\|^2
+
\frac{1}{2\eta}\|u-z\|^2
\right\}.
\]

On suppose qu'il existe une famille limite standardisée de problèmes, et qu'après normalisation des paramètres il existe un projecteur orthogonal déterministe \(P_0\) tel que, au point fixe pertinent \(u_\alpha^\star\),
\[
\|D\Pi_{M_\alpha}(u_\alpha^\star)-P_0\|_{\mathrm{op}}\xrightarrow[\alpha\to\infty]{}0
\]
uniformément sur la famille considérée.  
On suppose également que les coefficients convergent vers des limites déterministes :
\[
\lambda_{\mathrm{ctx},\alpha}\to \lambda_{\mathrm{ctx},\star},
\qquad
\lambda_{\mathrm{ridge},\alpha}\to \lambda_{\mathrm{ridge},\star}.
\]

\subsection{Hypothèse de concentration géométrique universelle}

On appelle \emph{hypothèse de concentration géométrique universelle} l'assertion suivante :

\begin{quote}
les hulls empiriques \(M_\alpha\) issus des embeddings de calibration, après normalisation de la représentation, possèdent un différentiel de projection qui se concentre autour d'un projecteur limite \(P_0\), uniformément sur la famille des corpus, encodeurs et PCA considérés.
\end{quote}

Cette hypothèse est le seul ingrédient non trivial nécessaire pour fermer le Cas 3 universel.

\subsection{Lemme 1 : Hessien de la distance carrée}

\begin{lemma}
Si \(\Pi_{M_\alpha}\) est différentiable en \(u\), alors
\[
\nabla^2 \dist(u,M_\alpha)^2
=
2\bigl(I-D\Pi_{M_\alpha}(u)\bigr).
\]
\end{lemma}

\begin{proof}
On utilise l'identité convexe classique
\[
\nabla\!\left(\frac12\dist(u,M_\alpha)^2\right)=u-\Pi_{M_\alpha}(u).
\]
En dérivant à nouveau,
\[
\nabla^2\!\left(\frac12\dist(u,M_\alpha)^2\right)
=
I-D\Pi_{M_\alpha}(u).
\]
En multipliant par \(2\), on obtient
\[
\nabla^2 \dist(u,M_\alpha)^2
=
2\bigl(I-D\Pi_{M_\alpha}(u)\bigr).
\]
\end{proof}

\subsection{Lemme 2 : concentration du Hessien géométrique}

\begin{lemma}
Sous l'hypothèse de concentration géométrique universelle,
\[
\nabla^2 \dist(u_\alpha^\star,M_\alpha)^2
\longrightarrow
2(I-P_0)
\]
en norme opératorielle.
\end{lemma}

\begin{proof}
Par le lemme précédent,
\[
\nabla^2 \dist(u_\alpha^\star,M_\alpha)^2
=
2\bigl(I-D\Pi_{M_\alpha}(u_\alpha^\star)\bigr).
\]
Donc
\[
\bigl\|
\nabla^2 \dist(u_\alpha^\star,M_\alpha)^2
-
2(I-P_0)
\bigr\|_{\mathrm{op}}
\le
2\|D\Pi_{M_\alpha}(u_\alpha^\star)-P_0\|_{\mathrm{op}}
\to 0.
\]
\end{proof}

\subsection{Théorème 1 : Hessien effectif universel}

\begin{theorem}
Sous l'hypothèse de concentration géométrique universelle,
\[
\nabla^2 J_\alpha(u_\alpha^\star;w,c)
\longrightarrow
H_\star
\]
avec
\[
H_\star
=
(1+2\lambda_{\mathrm{ridge},\star})I
+
2\lambda_{\mathrm{ctx},\star}(I-P_0).
\]
\end{theorem}

\begin{proof}
On a
\[
\nabla^2 J_\alpha(u_\alpha^\star;w,c)
=
I
+
2\lambda_{\mathrm{ridge},\alpha}I
+
\lambda_{\mathrm{ctx},\alpha}\nabla^2 \dist(u_\alpha^\star,M_\alpha)^2.
\]
En appliquant le lemme précédent,
\[
\nabla^2 J_\alpha(u_\alpha^\star;w,c)
=
I
+
2\lambda_{\mathrm{ridge},\alpha}I
+
2\lambda_{\mathrm{ctx},\alpha}\bigl(I-D\Pi_{M_\alpha}(u_\alpha^\star)\bigr).
\]
Puis la convergence des coefficients et de \(D\Pi_{M_\alpha}(u_\alpha^\star)\) donne
\[
\nabla^2 J_\alpha(u_\alpha^\star;w,c)\to
(1+2\lambda_{\mathrm{ridge},\star})I
+
2\lambda_{\mathrm{ctx},\star}(I-P_0)
=
H_\star.
\]
\end{proof}

\subsection{Théorème 2 : forme résolvante du Jacobien du prox}

\begin{theorem}
Sous les hypothèses précédentes,
\[
DT_{\eta,\alpha}(u_\alpha^\star;w,c)
\longrightarrow
(I+\eta H_\star)^{-1}
\]
en norme opératorielle.
\end{theorem}

\begin{proof}
Le point fixe \(u_\alpha^\star\) satisfait la condition d'optimalité du problème proximal. Dans la zone où la face active est stable, l'application résolvante est différentiable et le théorème des fonctions implicites donne
\[
DT_{\eta,\alpha}(u_\alpha^\star;w,c)
=
\bigl(I+\eta \nabla^2 J_\alpha(u_\alpha^\star;w,c)\bigr)^{-1}.
\]
Comme \(\nabla^2 J_\alpha(u_\alpha^\star;w,c)\to H_\star\), la continuité de l'inversion sur l'ensemble des matrices définies positives implique
\[
DT_{\eta,\alpha}(u_\alpha^\star;w,c)
\to
(I+\eta H_\star)^{-1}.
\]
\end{proof}

\subsection{Conséquence spectrale}

Comme \(P_0\) est un projecteur orthogonal, son spectre est contenu dans \(\{0,1\}\). Si \(r=\mathrm{rank}(P_0)\), alors
\[
\operatorname{Spec}(H_\star)
=
\left\{1+2\lambda_{\mathrm{ridge},\star}\right\}^{r}
\cup
\left\{1+2\lambda_{\mathrm{ridge},\star}+2\lambda_{\mathrm{ctx},\star}\right\}^{d-r}.
\]
Par conséquent,
\[
\operatorname{Spec}\bigl((I+\eta H_\star)^{-1}\bigr)
=
\left\{
\frac{1}{1+\eta(1+2\lambda_{\mathrm{ridge},\star})}
\right\}^{r}
\cup
\left\{
\frac{1}{1+\eta(1+2\lambda_{\mathrm{ridge},\star}+2\lambda_{\mathrm{ctx},\star})}
\right\}^{d-r}.
\]
Le rayon spectral local vaut donc
\[
\rho_\star
=
\frac{1}{1+\eta(1+2\lambda_{\mathrm{ridge},\star})}.
\]

\subsection{Cas 3 universel}

\begin{theorem}[Cas 3 universel]
Sous l'hypothèse de concentration géométrique universelle, et après normalisation de la famille de problèmes de sorte que les coefficients convergent vers des constantes déterministes \(\lambda_{\mathrm{ctx},\star}\) et \(\lambda_{\mathrm{ridge},\star}\), la dynamique locale du noyau PCH est gouvernée asymptotiquement par un opérateur résolvant universel :
\[
DT_{\eta,\alpha}(u_\alpha^\star;w,c)
\to
(I+\eta H_\star)^{-1},
\]
où
\[
H_\star
=
(1+2\lambda_{\mathrm{ridge},\star})I
+
2\lambda_{\mathrm{ctx},\star}(I-P_0),
\]
et où le rayon spectral local est
\[
\rho_\star
=
\frac{1}{1+\eta(1+2\lambda_{\mathrm{ridge},\star})}.
\]
\end{theorem}

\begin{proof}
C'est la combinaison directe du théorème sur le Hessien effectif et du théorème sur le Jacobien du prox.
\end{proof}

\subsection{Lecture finale}

Le Cas 3 universel ne repose pas sur une magie du prox, mais sur une chaîne logique unique :

\[
\text{concentration géométrique du hull}
\;\Longrightarrow\;
\text{stabilité de la face active}
\;\Longrightarrow\;
D\Pi_{M_\alpha}(u_\alpha^\star)\to P_0
\;\Longrightarrow\;
\nabla^2 J_\alpha(u_\alpha^\star)\to H_\star
\;\Longrightarrow\;
DT_{\eta,\alpha}(u_\alpha^\star)\to (I+\eta H_\star)^{-1}.
\]

La seule hypothèse réellement substantielle est donc la concentration géométrique uniforme du différentiel de projection autour d'un projecteur limite \(P_0\).
Une fois cette hypothèse admise, le Cas 3 universel devient une conséquence analytique.




%%%%%%%%%%%%%%%%%%%%%%%%%%%%%%%%%%%%%%%%%%%%%%%%%%%%%%%%%%%%%%%%%%%%%%%%%%%%%%%%%%%%%%%%%%%%%%%%%%%%%%%%%%%%%%%%%%%%%%%%%%%%%%%%%%%%%%%%%%%%%%%%%%%%%%%%%%%%%%%%%%%%%%%%%%%%%%%%%%%%%%%%%%%%%%%%%%%%%%%%%%%%%%%%%%%%%%%%%%%%%%%%%%%%%%%%%%%%%%%%%%%%%%%%%%%%%%%%%%%%%%%%%%%%%%%%%%%%%%%%%%%impossibilité d universalité

\section{Impossibilité d'une universalité absolue indépendante des hyperparamètres}

\subsection{Cadre}

On considère une famille de problèmes PCH indexée par \(\alpha\), avec énergie locale
\[
J_\alpha(u;w,c)
=
\frac12\|u-w\|^2
+
\lambda_{\mathrm{ctx},\alpha}\,\dist(u,M_\alpha)^2
+
\lambda_{\mathrm{ridge},\alpha}\,\|u\|^2,
\]
et point fixe local pertinent \(u_\alpha^\star\).

On suppose que la projection est différentiable en \(u_\alpha^\star\), de sorte que
\[
\nabla^2 \dist(u_\alpha^\star,M_\alpha)^2
=
2\bigl(I-D\Pi_{M_\alpha}(u_\alpha^\star)\bigr).
\]

On note
\[
P_\alpha := D\Pi_{M_\alpha}(u_\alpha^\star).
\]

\subsection{Lemme : forme exacte du Hessien local}

Sous différentiabilité de \(\Pi_{M_\alpha}\) en \(u_\alpha^\star\), on a
\[
\nabla^2 J_\alpha(u_\alpha^\star;w,c)
=
I
+
2\lambda_{\mathrm{ridge},\alpha}I
+
2\lambda_{\mathrm{ctx},\alpha}\bigl(I-P_\alpha\bigr).
\]

\begin{proof}
On dérive terme à terme :
\[
\nabla^2\Big(\frac12\|u-w\|^2\Big)=I,
\qquad
\nabla^2\big(\lambda_{\mathrm{ridge},\alpha}\|u\|^2\big)=2\lambda_{\mathrm{ridge},\alpha}I,
\]
et, par la formule de Hessien de la distance carrée,
\[
\nabla^2\big(\lambda_{\mathrm{ctx},\alpha}\dist(u,M_\alpha)^2\big)
=
2\lambda_{\mathrm{ctx},\alpha}(I-P_\alpha).
\]
En sommant, on obtient la formule annoncée.
\end{proof}

\subsection{Lemme : Jacobien du proximal}

Supposons que le proximal local soit différentiable au point fixe. Alors
\[
DT_{\eta,\alpha}(u_\alpha^\star;w,c)
=
\Bigl(I+\eta \nabla^2 J_\alpha(u_\alpha^\star;w,c)\Bigr)^{-1}.
\]
Par conséquent,
\[
DT_{\eta,\alpha}(u_\alpha^\star;w,c)
=
\Bigl(
I+\eta\bigl[
I+2\lambda_{\mathrm{ridge},\alpha}I
+2\lambda_{\mathrm{ctx},\alpha}(I-P_\alpha)
\bigr]
\Bigr)^{-1}.
\]

\begin{proof}
C'est l'identité standard du résolvant d'un problème fortement convexe localement \(C^2\), obtenue par le théorème des fonctions implicites appliqué à la condition d'optimalité du proximal.
\end{proof}

\subsection{Proposition : dépendance nécessaire aux hyperparamètres}

Le Jacobien local dépend nécessairement de \(\lambda_{\mathrm{ridge},\alpha}\) et de \(\lambda_{\mathrm{ctx},\alpha}\), sauf dégénérescence.

\begin{proof}
Considérons d'abord la restriction de \(DT_{\eta,\alpha}(u_\alpha^\star;w,c)\) à une direction tangentielle \(v\) telle que \(P_\alpha v=v\). Alors
\[
(I-P_\alpha)v=0,
\]
donc
\[
DT_{\eta,\alpha}(u_\alpha^\star;w,c)v
=
\frac{1}{1+\eta(1+2\lambda_{\mathrm{ridge},\alpha})}v.
\]
Ainsi, le facteur spectral tangent dépend explicitement de \(\lambda_{\mathrm{ridge},\alpha}\).

De même, sur une direction normale \(n\) telle que \(P_\alpha n=0\), on a
\[
(I-P_\alpha)n=n,
\]
d'où
\[
DT_{\eta,\alpha}(u_\alpha^\star;w,c)n
=
\frac{1}{1+\eta(1+2\lambda_{\mathrm{ridge},\alpha}+2\lambda_{\mathrm{ctx},\alpha})}n.
\]
Le facteur spectral normal dépend explicitement de \(\lambda_{\mathrm{ridge},\alpha}\) et de \(\lambda_{\mathrm{ctx},\alpha}\).

Il s'ensuit que, sauf si ces coefficients sont identiques pour toute la famille ou si l'on effectue une normalisation compensatrice explicite, le Jacobien local ne peut pas être indépendant des hyperparamètres.
\end{proof}

\subsection{Théorème : impossibilité d'une universalité absolue sans normalisation}

Supposons qu'il existe une matrice fixe \(L\) telle que, pour toute famille indexée par \(\alpha\),
\[
DT_{\eta,\alpha}(u_\alpha^\star;w,c)\equiv L
\]
indépendamment de \(\lambda_{\mathrm{ridge},\alpha}\) et \(\lambda_{\mathrm{ctx},\alpha}\).

Alors, nécessairement, les coefficients \(\lambda_{\mathrm{ridge},\alpha}\) et \(\lambda_{\mathrm{ctx},\alpha}\) doivent être constants sur la famille, ou bien la famille doit être renormalisée de sorte que les variations de ces coefficients soient absorbées.

\begin{proof}
Supposons par l'absurde qu'il existe deux indices \(\alpha\) et \(\beta\) tels que
\[
\lambda_{\mathrm{ridge},\alpha}\neq \lambda_{\mathrm{ridge},\beta}.
\]
Choisissons une direction tangentielle \(v\) commune aux deux configurations lorsque cela a un sens local, ou, à défaut, considérons les valeurs propres tangentielles du Jacobien local. On obtient alors
\[
\rho_{\alpha}^{T}
=
\frac{1}{1+\eta(1+2\lambda_{\mathrm{ridge},\alpha})},
\qquad
\rho_{\beta}^{T}
=
\frac{1}{1+\eta(1+2\lambda_{\mathrm{ridge},\beta})}.
\]
Comme la fonction
\[
x\mapsto \frac{1}{1+\eta(1+2x)}
\]
est strictement monotone sur \(\mathbb R_+\), on a
\[
\rho_{\alpha}^{T}\neq \rho_{\beta}^{T}.
\]
Donc les Jacobiens locaux ne peuvent pas coïncider.

Le même raisonnement vaut pour \(\lambda_{\mathrm{ctx},\alpha}\) via les valeurs propres normales.

Ainsi, une universalité absolue indépendante des hyperparamètres est impossible dans ce modèle, sauf si l'on impose que les coefficients ne varient pas ou que leur variation soit compensée par une normalisation externe.
\end{proof}

\subsection{Corollaire : forme correcte du Cas 3}

La forme mathématiquement défendable du Cas 3 est donc la suivante :

\[
DT_{\eta,\alpha}(u_\alpha^\star;w,c)
\to
(I+\eta H_{\star,\alpha})^{-1},
\]
avec
\[
H_{\star,\alpha}
=
(1+2\lambda_{\mathrm{ridge},\alpha})I
+
2\lambda_{\mathrm{ctx},\alpha}(I-P_0),
\]
où \(P_0\) est le projecteur limite géométrique.

Autrement dit, il existe une universalité \emph{géométrique} via \(P_0\), mais pas une universalité \emph{absolue} des coefficients sans hypothèse supplémentaire de normalisation.

\subsection{Conclusion}

Le modèle PCH impose mathématiquement que la forme du Jacobien local dépende des coefficients de régularisation, sauf si l'on fixe ces coefficients ou si l'on travaille sur une famille normalisée.

Ainsi :

\begin{itemize}
    \item l'universalité \emph{absolue} indépendante des hyperparamètres est impossible dans le modèle brut ;
    \item l'universalité \emph{conditionnelle} après normalisation est possible ;
    \item la bonne cible théorique du Cas 3 est donc une universalité géométrique du projecteur limite \(P_0\), non une invariance brute des paramètres.
\end{itemize}





%%%%%%%%%%%%%%%%%%%%%%%%%%%%%%%%%%%%%%%%%%%%%%%%%%%%%%%%%%%%%%%%%%%%%%%%%%%%%%%%%%%%%%%%%%%%%%%%%%%%%%%%%%%%%%%%%%%%%%%%%%%%%%%%%%%%%%%%%%%%%%%%%%%%%%%%%%%%%%%%%%%%%%%%%%%%%%%%%%%%%%%%%%%%%%%%%%%%%%%%%%%%%%%%%%%%%%%%%%%%%%%%%%%%%%%%%%%%%%%%%%%%%%%%%%%%%%%%%%%%%%%%%%%%%%%%%%%%%%%%%%%Cas 3 universel sous normalisation de courbure

\section{Cas 3 universel sous normalisation de courbure}

\subsection{Cadre}

On considère une famille de problèmes PCH indexée par \(\alpha\), où \(\alpha\) code le corpus, l'encodeur, la PCA, la dimension latente et les hyperparamètres.

Pour chaque \(\alpha\), on a un ensemble admissible convexe fermé
\[
M_\alpha \subset \mathbb R^{d_\alpha},
\]
et une énergie locale
\[
J_\alpha(u;w,c)
=
\frac12\|u-w\|^2
+
\lambda_{\mathrm{ctx},\alpha}\,\dist(u,M_\alpha)^2
+
\lambda_{\mathrm{ridge},\alpha}\,\|u\|^2.
\]

On note
\[
u_\alpha^\star
\]
le point fixe local pertinent, et
\[
P_\alpha := D\Pi_{M_\alpha}(u_\alpha^\star)
\]
le différentiel de la projection convexe au point fixe.

On introduit également
\[
\mu_\alpha := 1+2\lambda_{\mathrm{ridge},\alpha},
\]
qui est la courbure tangentielle de base induite par le terme quadratique radial.

\subsection{Hypothèse de normalisation universelle}

La normalisation choisie dans le contexte de la théorie est la suivante.

\begin{enumerate}
    \item \textbf{Normalisation de courbure tangentielle.}
    On travaille avec la quantité adimensionnée
    \[
    \widehat J_\alpha := \frac{1}{\mu_\alpha}J_\alpha.
    \]
    Cette normalisation fixe l'échelle tangentielle à l'unité.

    \item \textbf{Normalisation du poids contextuel.}
    On suppose qu'il existe une constante universelle \(\gamma>0\) telle que
    \[
    \gamma_\alpha := \frac{\lambda_{\mathrm{ctx},\alpha}}{\mu_\alpha} \equiv \gamma
    \quad\text{pour toute } \alpha.
    \]

    \item \textbf{Normalisation du pas proximal.}
    On suppose qu'il existe une constante universelle \(\tau>0\) telle que
    \[
    \tau_\alpha := \eta_\alpha \mu_\alpha \equiv \tau
    \quad\text{pour toute } \alpha.
    \]

    \item \textbf{Concentration géométrique universelle.}
    On suppose qu'il existe un projecteur orthogonal déterministe \(P_0\) tel que
    \[
    \|P_\alpha - P_0\|_{\mathrm{op}} \to 0
    \qquad\text{uniformément sur la famille considérée.}
    \]
\end{enumerate}

Cette hypothèse est la bonne forme de normalisation pour une théorie universelle : elle fixe l'unité locale de courbure, puis elle impose que le rapport contextuel et le pas proximal soient mesurés dans cette unité.

\subsection{Hessien local normalisé}

Sous différentiabilité de la projection en \(u_\alpha^\star\), on a l'identité classique
\[
\nabla^2 \dist(u_\alpha^\star,M_\alpha)^2
=
2\bigl(I-P_\alpha\bigr).
\]

Donc
\[
\nabla^2 J_\alpha(u_\alpha^\star;w,c)
=
\mu_\alpha I
+
2\lambda_{\mathrm{ctx},\alpha}(I-P_\alpha).
\]

En divisant par \(\mu_\alpha\), le Hessien normalisé vaut
\[
\widehat H_\alpha
:=
\frac{1}{\mu_\alpha}\nabla^2 J_\alpha(u_\alpha^\star;w,c)
=
I + 2\gamma (I-P_\alpha).
\]

Par conséquent,
\[
\widehat H_\alpha \to \widehat H_\star
:=
I + 2\gamma(I-P_0)
\]
en norme opératorielle.

\subsection{Jacien du proximal normalisé}

Le proximal local vérifie, au point fixe,
\[
DT_{\eta_\alpha,\alpha}(u_\alpha^\star;w,c)
=
\bigl(I+\eta_\alpha \nabla^2 J_\alpha(u_\alpha^\star;w,c)\bigr)^{-1}.
\]

Comme \(\eta_\alpha\mu_\alpha=\tau\), on obtient
\[
DT_{\eta_\alpha,\alpha}(u_\alpha^\star;w,c)
=
\bigl(I+\tau \widehat H_\alpha\bigr)^{-1}.
\]

Ainsi,
\[
DT_{\eta_\alpha,\alpha}(u_\alpha^\star;w,c)
\to
\bigl(I+\tau \widehat H_\star\bigr)^{-1}
=
\bigl(I+\tau[I+2\gamma(I-P_0)]\bigr)^{-1}.
\]

La convergence est en norme opératorielle, car l'application \(A\mapsto A^{-1}\) est continue sur l'ensemble des matrices uniformément inversibles et, ici, les matrices
\[
I+\tau \widehat H_\alpha
\]
sont uniformément définies positives.

\subsection{Structure spectrale universelle}

Comme \(P_0\) est un projecteur orthogonal, son spectre est contenu dans \(\{0,1\}\).

Si \(r=\operatorname{rank}(P_0)\), alors sur l'espace tangent \(\operatorname{Ran}(P_0)\),
\[
(I-P_0)v=0,
\]
et sur l'espace normal \(\ker(P_0)\),
\[
(I-P_0)n=n.
\]

Donc les valeurs propres de \(\widehat H_\star\) sont :
\[
1 \quad\text{sur les directions tangentes},
\qquad
1+2\gamma \quad\text{sur les directions normales}.
\]

Par conséquent,
\[
\operatorname{Spec}\bigl((I+\tau \widehat H_\star)^{-1}\bigr)
=
\left\{
\frac{1}{1+\tau}
\right\}^{r}
\cup
\left\{
\frac{1}{1+\tau(1+2\gamma)}
\right\}^{d_\alpha-r}.
\]

Le rayon spectral universel local est donc
\[
\rho_\star = \frac{1}{1+\tau}.
\]

\subsection{Théorème de Cas 3 universel}

\begin{theorem}[Cas 3 universel sous normalisation de courbure]
Sous les hypothèses de normalisation universelle ci-dessus,
la dynamique locale du noyau PCH est gouvernée par un opérateur résolvant universel :
\[
DT_{\eta_\alpha,\alpha}(u_\alpha^\star;w,c)
\to
\bigl(I+\tau[I+2\gamma(I-P_0)]\bigr)^{-1},
\]
indépendamment du corpus, de l'encodeur, de la PCA et des hyperparamètres originaux,
dès lors que la normalisation par \(\mu_\alpha\) est imposée.

En particulier, le spectre local limite est universel et vaut
\[
\operatorname{Spec}
=
\left\{
\frac{1}{1+\tau}
\right\}^{r}
\cup
\left\{
\frac{1}{1+\tau(1+2\gamma)}
\right\}^{d_\alpha-r}.
\]
\end{theorem}

\begin{proof}
On part de l'identité exacte
\[
DT_{\eta_\alpha,\alpha}(u_\alpha^\star;w,c)
=
\bigl(I+\eta_\alpha \nabla^2 J_\alpha(u_\alpha^\star;w,c)\bigr)^{-1}.
\]
En écrivant
\[
\nabla^2 J_\alpha(u_\alpha^\star;w,c)
=
\mu_\alpha \widehat H_\alpha,
\]
on obtient
\[
DT_{\eta_\alpha,\alpha}(u_\alpha^\star;w,c)
=
\bigl(I+\eta_\alpha\mu_\alpha \widehat H_\alpha\bigr)^{-1}
=
\bigl(I+\tau \widehat H_\alpha\bigr)^{-1}.
\]
Par l'hypothèse de concentration géométrique universelle,
\[
\widehat H_\alpha \to \widehat H_\star = I+2\gamma(I-P_0),
\]
donc, par continuité de l'inversion sur les matrices uniformément inversibles,
\[
DT_{\eta_\alpha,\alpha}(u_\alpha^\star;w,c)
\to
\bigl(I+\tau \widehat H_\star\bigr)^{-1}.
\]
Le calcul spectral de \(\widehat H_\star\) donne immédiatement la forme annoncée.
\end{proof}

\subsection{Corollaire : version universelle du Hessien effectif}

Sous la même normalisation, le Hessien effectif universel est
\[
H_{\mathrm{eff}}^\star = I+2\gamma(I-P_0),
\]
et non pas une matrice dépendant de \(\alpha\).

Autrement dit, toute la dépendance aux corpus, encodeurs, PCA et hyperparamètres a été absorbée par la jauge
\[
\mu_\alpha=1+2\lambda_{\mathrm{ridge},\alpha}
\]
et par la normalisation des ratios
\[
\gamma_\alpha=\frac{\lambda_{\mathrm{ctx},\alpha}}{\mu_\alpha},
\qquad
\tau_\alpha=\eta_\alpha\mu_\alpha.
\]

\subsection{Lecture finale}

La bonne lecture du Cas 3 universel est la suivante :

\begin{itemize}
    \item la géométrie aléatoire des hulls se concentre vers un projecteur limite \(P_0\) ;
    \item la normalisation de courbure fixe l'échelle tangentielle ;
    \item le rapport contextuel et le pas proximal sont exprimés dans cette jauge ;
    \item le Jacobien du proximal converge alors vers un opérateur résolvant universel ;
    \item la loi spectrale limite devient indépendante de l'origine des données.
\end{itemize}

Ainsi, le Cas 3 universel n'est pas l'affirmation naïve que tous les hyperparamètres disparaissent sans hypothèse : c'est le théorème que, \emph{dans la bonne jauge}, la dynamique locale du noyau PCH devient universelle.






%%%%%%%%%%%%%%%%%%%%%%%%%%%%%%%%%%%%%%%%%%%%%%%%%%%%%%%%%%%%%%%%%%%%%%%%%%%%%%%%%%%%%%%%%%%%%%%%%%%%%%%%%%%%%%%%%%%%%%%%%%%%%%%%%%%%%%%%%%%%%%%%%%%%%%%%%%%%%%%%%%%%%%%%%%%%%%%%%%%%%%%%%%%%%%%%%%%%%%%%%%%%%%%%%%%%%%%%%%%%%%%%%%%%%%%%%%%%%%%%%%%%%%%%%%%%%%%%%%%%%%%%%%%%%%%%%%%%%%%%%%%raport de l'experience cas 3 universel normalisé  pch_cas3_normalized_universal.py pour la version de "Cas 3 universel sous normalisation de courbure"(sans Réduction des hypothèses H2 et H3 à des axiomes plus fondamentaux)

\section{Rapport rigoureux sur le Cas 3 normalisé et ses trois campagnes expérimentales}

\subsection{Objet et jauge}

On étudie la dynamique locale du noyau PCH sous la jauge normalisée
\[
\mu = 1 + 2\lambda_{\mathrm{ridge}}, 
\qquad
\gamma = \frac{\lambda_{\mathrm{ctx}}}{\mu},
\qquad
\tau = \eta \mu,
\]
avec \(\gamma_{\mathrm{target}}=0.5\) et \(\tau_{\mathrm{target}}=3.4\) dans toutes les campagnes.

La valeur théorique scalaire associée est donc
\[
\rho_\star = \frac{1}{1+\tau} = \frac{1}{4.4} = 0.22727272727272727.
\]

L’objectif expérimental est de vérifier, dans cette jauge, si la dynamique locale du prox PCH se comporte comme un opérateur résolvant normalisé
\[
DT_{\eta,\alpha}^{\mathrm{PCH}}(u_\alpha^\star;w,c)
\approx
\bigl(I+\tau[I+2\gamma(I-P_\alpha)]\bigr)^{-1},
\]
et si un projecteur limite \(P_0\) peut être mis en évidence à travers plusieurs corpus, encodeurs et partitions.

\subsection{Protocole exact observé dans les sorties}

Le script fixe explicitement :
\[
\texttt{pca\_fit}=\texttt{calibration\_only},
\qquad
\texttt{context\_definition}=\texttt{convex hull of calibration embeddings in latent space},
\]
et utilise
\[
\texttt{exact\_projection}=\texttt{True},\quad
\texttt{exact\_pch}=\texttt{True},\quad
\texttt{exact\_proximal\_map}=\texttt{True},\quad
\texttt{local\_linearization}=\texttt{finite\_difference\_jacobian}.
\]

Les sorties collées fixent de manière identique dans toutes les campagnes :
\[
\texttt{latent\_dim}=48,\qquad
n_{\text{train}}=128,\qquad
n_{\text{test}}=128,\qquad
n_{\text{calib}}=48 \text{ par contexte},
\]
\[
n_{\text{probe train}}=n_{\text{probe test}}=32 \text{ par contexte},\qquad
\texttt{max\_iters}=50,\qquad
\texttt{seed}=7,\qquad
\texttt{batch\_size}=64,
\]
\[
\texttt{energy\_check\_samples}=4,\qquad
\texttt{jacobian\_samples}=4,\qquad
\texttt{stop\_tol}=10^{-11},\qquad
\texttt{useful\_ratio\_relative\_floor}=10^{-9}.
\]

Le hull de chaque contexte est construit comme
\[
M_c = \operatorname{conv}(D_c),
\]
où \(D_c\) contient exactement \(48\) embeddings de calibration du contexte \(c\).

Les sorties montrent en outre que, pour chaque valeur de \(\lambda_{\mathrm{ridge}}\), les runs \texttt{regularizer\_kind = ridge} et \texttt{regularizer\_kind = none} sont numériquement identiques à la précision reportée ; les tableaux ci-dessous ne gardent donc qu’une ligne par valeur de \(\lambda_{\mathrm{ridge}}\).

Le lien entre les hyperparamètres bruts et la jauge normalisée est explicite :
\[
\lambda_{\mathrm{ctx,raw}}=\gamma \mu,\qquad
\eta_{\mathrm{raw}}=\frac{\tau}{\mu}.
\]
Ainsi,
\[
\lambda_{\mathrm{ridge}}\in\{0.2,\,0.8,\,1.5,\,3.0\}
\quad\Longrightarrow\quad
\mu\in\{1.4,\,2.6,\,4.0,\,7.0\},
\]
et donc
\[
\lambda_{\mathrm{ctx,raw}}\in\{0.7,\,1.3,\,2.0,\,3.5\},
\qquad
\eta_{\mathrm{raw}}\in\left\{2.428571428571429,\ 1.3076923076923077,\ 0.85,\ 0.4857142857142857\right\}.
\]



\subsection{Définitions des diagnostics}

Le script calcule les diagnostics numériques suivants, tous tels qu’ils sont définis dans le code expérimental.

On note d’abord
\[
R_{\mathrm{pred}} = \bigl(I+\tau[I+2\gamma(I-P_\alpha)]\bigr)^{-1}.
\]
Puis :
\[
\texttt{rho\_obs}=\|DT_{\hat{}}\|_2,\qquad
\texttt{rho\_pred}=\|R_{\mathrm{pred}}\|_2,
\qquad
\texttt{rho\_gap}=|\texttt{rho\_obs}-\texttt{rho\_pred}|.
\]

Les autres écarts sont définis par la norme opératorielle spectrale (\(\|\cdot\|_2\)) :
\[
\texttt{operator\_gap}=\|DT_{\hat{}}-R_{\mathrm{pred}}\|_2,
\]
\[
\texttt{projector\_gap}=\|P_{\hat{}}^2-P_{\hat{}}\|_2,
\]
\[
\texttt{projector\_symmetry\_gap}=\|P_{\hat{}}-P_{\hat{}}^\top\|_2.
\]

Enfin,
\[
\texttt{projector\_trace}=\operatorname{tr}(P_{\hat{}}).
\]

Autrement dit, \texttt{operator\_gap}, \texttt{projector\_gap} et \texttt{projector\_symmetry\_gap} sont tous des écarts mesurés en norme spectrale, et non en norme de Frobenius.

\subsection{Validation quantitative de la loi normalisée}

Les sorties reportent directement, pour chaque run, les diagnostics permettant de comparer l’opérateur observé à l’opérateur prédit :
\[
R_{\mathrm{pred}} = \bigl(I+\tau[I+2\gamma(I-P_\alpha)]\bigr)^{-1}.
\]
Les diagnostics sont :
\[
\texttt{rho\_obs},\quad \texttt{rho\_pred},\quad
\texttt{rho\_gap},\quad
\texttt{operator\_gap},\quad
\texttt{projector\_gap},\quad
\texttt{projector\_symmetry\_gap},\quad
\texttt{projector\_trace}.
\]
La loi scalaire \(\rho_\star=1/(1+\tau)\) est fortement corroborée, mais les écarts opérateur/projecteur ne sont pas nuls uniformément sur toute la campagne ; ils sont précisément ceux que le fichier de résultats consigne.

\subsection{Résumé global des écarts sur \(\rho\) et des traces}

Les statistiques globales suivantes condensent l’ensemble des runs collés :

\begin{center}
\small
\begin{tabular}{lccc}
\hline
Campagne & \(\max |\rho_{\mathrm{obs}}-\rho_\star|\) & Trace train (min--max) & Trace test (min--max) \\
\hline
ex1 & \(3.2508604025371923\times 10^{-6}\) & \(15.718750000000124\)--\(21.937569631649836\) & \(15.521878804368042\)--\(21.929687499999694\) \\
ex2 & \(8.993255675715517\times 10^{-7}\) & \(15.167591529956862\)--\(21.726562500000135\) & \(15.557274580859561\)--\(21.703197696365727\) \\
ex3 & \(1.4340065024320836\times 10^{-8}\) & \(15.468750000000911\)--\(22.687499999998956\) & \(15.479688630394707\)--\(22.36718749999978\) \\
\hline
\end{tabular}
\end{center}

Sur l’ensemble de tous les runs collés,
\[
\max_{\text{tous runs}}|\rho_{\mathrm{obs}}-\rho_\star|
=
3.2508604025371923\times 10^{-6}.
\]

Les traces de projecteur observées sur l’ensemble complet des runs couvrent
\[
\operatorname{tr}(P_\alpha)\in
[15.167591529956862,\ 22.687499999998956]
\]
sur train, et
\[
\operatorname{tr}(P_\alpha)\in
[15.479688630394707,\ 22.36718749999978]
\]
sur test.

\subsection{Campagne ex1 : NEWS4-A, encodeur \texttt{sentence-transformers/all-MiniLM-L6-v2}}

Cette campagne utilise
\[
\text{NEWS4-A}=
\{\texttt{alt.atheism},\texttt{comp.graphics},\texttt{sci.space},\texttt{talk.politics.misc}\}.
\]
Le probe brut vaut, dans tous les runs de cette campagne,
\[
\texttt{probe\_raw.accuracy}=0.75,
\qquad
\texttt{probe\_raw.macro\_f1}=0.751195474136017.
\]

\begin{center}
\small
\begin{tabular}{c c c c c c c c}
\hline
\(\lambda_{\mathrm{ridge}}\) & \(\mu\) & \(\rho_{\mathrm{obs}}\) & \(\rho_{\mathrm{pred}}\) & \(\texttt{rho\_gap}\) & \(\texttt{op\_gap}\) & \(\texttt{proj\_gap}\) & \(\texttt{sym\_gap}\) \\
\hline
0.2 & 1.4 & 0.2272727272773883 & 0.22728282155581642 & 1.0094286038568687e-05 & 0.0001558958667545639 & 0.0008670918503703889 & 0.0007320302848481992 \\
0.8 & 2.6 & 0.2272727272784203 & 0.22728249658544594 & 9.769316633252679e-06 & 0.00011218023075833402 & 0.0005873142297521498 & 0.0005176504744980861 \\
1.5 & 4.0 & 0.2272759781331298 & 0.2272730382308339 & 2.939902295909843e-06 & 0.000134740977098282 & 0.0009651409107477826 & 0.00016929495456717605 \\
3.0 & 7.0 & 0.2272727272796613 & 0.22727272727371062 & 5.9506908949014115e-12 & 6.127152827634715e-06 & 5.863884998453531e-05 & 7.049770118443349e-05 \\
\hline
\end{tabular}
\end{center}

\begin{center}
\small
\begin{tabular}{c c c c c c c c}
\hline
\(\lambda_{\mathrm{ridge}}\) & tr train & tr test & acc brute & F1 brute & acc oracle & F1 oracle & acc routée / ctx train / ctx test \\
\hline
0.2 & 15.718750000000124 & 15.521878804368042 & 0.75 & 0.751195474136017 & 0.8984375 & 0.898663565372735 & 0.8125 / 0.84375 / 0.8359375 \\
0.8 & 20.679612487094683 & 21.014599993049256 & 0.75 & 0.751195474136017 & 0.9765625 & 0.9765512265512266 & 0.8046875 / 0.84375 / 0.8359375 \\
1.5 & 21.84374999999988 & 21.746580897633194 & 0.75 & 0.751195474136017 & 1.0 & 1.0 & 0.796875 / 0.84375 / 0.8359375 \\
3.0 & 21.937569631649836 & 21.929687499999694 & 0.75 & 0.751195474136017 & 1.0 & 1.0 & 0.8203125 / 0.8828125 / 0.8359375 \\
\hline
\end{tabular}
\end{center}

Les gains exacts d’accuracy par rapport au probe brut sont donc :
\[
\Delta_{\mathrm{oracle}} \in \{0.1484375,\ 0.2265625,\ 0.25,\ 0.25\},
\]
\[
\Delta_{\mathrm{routed}} \in \{0.0625,\ 0.0546875,\ 0.046875,\ 0.0703125\}.
\]

La monotonie énergétique est parfaite :
\[
\texttt{energy\_monotonic\_violations}=0
\]
dans tous les runs.

Les résidus de point fixe restent au niveau machine, sur train et test, entre
\[
5.719551864242296\times 10^{-17}
\quad\text{et}\quad
8.233436468072948\times 10^{-17}.
\]

Les traces de projecteur augmentent nettement avec \(\lambda_{\mathrm{ridge}}\), passant globalement d’environ \(15.7\) à environ \(21.9\), ce qui montre que la géométrie active varie avec la jauge même lorsque la loi scalaire reste très proche de sa valeur théorique.

\subsection{Campagne ex2 : NEWS4-A, encodeur \texttt{sentence-transformers/all-mpnet-base-v2}}

Cette campagne utilise le même split de classes que ex1, mais avec un autre encodeur :
\[
\texttt{sentence-transformers/all-mpnet-base-v2}.
\]
Le probe brut vaut, dans tous les runs de cette campagne,
\[
\texttt{probe\_raw.accuracy}=0.765625,
\qquad
\texttt{probe\_raw.macro\_f1}=0.7655654909753271.
\]

\begin{center}
\small
\begin{tabular}{c c c c c c c c}
\hline
\(\lambda_{\mathrm{ridge}}\) & \(\mu\) & \(\rho_{\mathrm{obs}}\) & \(\rho_{\mathrm{pred}}\) & \(\texttt{rho\_gap}\) & \(\texttt{op\_gap}\) & \(\texttt{proj\_gap}\) & \(\texttt{sym\_gap}\) \\
\hline
0.2 & 1.4 & 0.22727362659829484 & 0.2272738558111333 & 1.7186317004824498e-06 & 0.0001554931757878567 & 0.0018626603742000045 & 0.0005518728469858196 \\
0.8 & 2.6 & 0.22727272727870296 & 0.22727272727354847 & 5.154480791844107e-12 & 1.50882290824329e-11 & 8.939010993638839e-12 & 1.4082735874628381e-11 \\
1.5 & 4.0 & 0.22727272727834164 & 0.22727272727354159 & 4.800021599811466e-12 & 1.3493074063578146e-11 & 8.550902904931793e-12 & 1.3535790047090683e-11 \\
3.0 & 7.0 & 0.2272727272796613 & 0.22727272727371062 & 5.9506908949014115e-12 & 1.6525333577874244e-11 & 9.75392014487398e-12 & 1.5373049632798975e-11 \\
\hline
\end{tabular}
\end{center}

\begin{center}
\small
\begin{tabular}{c c c c c c c c}
\hline
\(\lambda_{\mathrm{ridge}}\) & tr train & tr test & acc brute & F1 brute & acc oracle & F1 oracle & acc routée / ctx train / ctx test \\
\hline
0.2 & 15.167591529956862 & 15.557274580859561 & 0.765625 & 0.7655654909753271 & 0.890625 & 0.8904558982683983 & 0.7578125 / 0.8828125 / 0.828125 \\
0.8 & 20.13281250000063 & 20.31250000000035 & 0.765625 & 0.7655654909753271 & 0.96875 & 0.9687072490351178 & 0.765625 / 0.8828125 / 0.828125 \\
1.5 & 21.390624999999755 & 21.52343749999962 & 0.765625 & 0.7655654909753271 & 0.9921875 & 0.9921855921855922 & 0.765625 / 0.8828125 / 0.828125 \\
3.0 & 21.726562500000135 & 21.703197696365727 & 0.765625 & 0.7655654909753271 & 1.0 & 1.0 & 0.78125 / 0.8828125 / 0.828125 \\
\hline
\end{tabular}
\end{center}

Les gains exacts sont donc :
\[
\Delta_{\mathrm{oracle}} \in \{0.125,\ 0.203125,\ 0.2265625,\ 0.234375\},
\]
\[
\Delta_{\mathrm{routed}} \in \{-0.0078125,\ 0,\ 0,\ 0.015625\}.
\]

La monotonie énergétique est encore parfaite :
\[
\texttt{energy\_monotonic\_violations}=0
\]
dans tous les runs.

Les résidus de point fixe restent au niveau machine, entre
\[
7.518524274569321\times 10^{-17}
\quad\text{et}\quad
8.10967773180263\times 10^{-17}.
\]

Les traces de projecteur passent globalement d’environ \(15.2\) à environ \(21.7\) lorsque \(\lambda_{\mathrm{ridge}}\) augmente. La loi scalaire normalisée reste cependant très stable, avec une séparation nette entre la partie locale résolvante et la partie aval de classification.

\subsection{Campagne ex3 : NEWS4-B, encodeur \texttt{sentence-transformers/all-MiniLM-L6-v2}}

Cette campagne utilise une autre partition du corpus :
\[
\text{NEWS4-B}=
\{\texttt{rec.motorcycles},\texttt{comp.os.ms-windows.misc},\texttt{sci.electronics},\texttt{sci.med}\}.
\]
Le probe brut vaut, dans tous les runs de cette campagne,
\[
\texttt{probe\_raw.accuracy}=0.7734375,
\qquad
\texttt{probe\_raw.macro\_f1}=0.7675757575757576.
\]

\begin{center}
\small
\begin{tabular}{c c c c c c c c}
\hline
\(\lambda_{\mathrm{ridge}}\) & \(\mu\) & \(\rho_{\mathrm{obs}}\) & \(\rho_{\mathrm{pred}}\) & \(\texttt{rho\_gap}\) & \(\texttt{op\_gap}\) & \(\texttt{proj\_gap}\) & \(\texttt{sym\_gap}\) \\
\hline
0.2 & 1.4 & 0.2272727416127923 & 0.22727760018715126 & 4.858581945244109e-06 & 0.00014124697270368887 & 0.0012600322188579549 & 0.0006089206336516322 \\
0.8 & 2.6 & 0.2272727272784203 & 0.22728249658544594 & 9.769316633252679e-06 & 0.00011218023075833402 & 0.0005873142297521498 & 0.0005176504744980861 \\
1.5 & 4.0 & 0.22727272727834164 & 0.22727272727354159 & 4.800021599811466e-12 & 1.3493074063578146e-11 & 8.550902904931793e-12 & 1.3535790047090683e-11 \\
3.0 & 7.0 & 0.2272727272796613 & 0.22727272727371062 & 5.9506908949014115e-12 & 1.6525333577874244e-11 & 9.75392014487398e-12 & 1.5373049632798975e-11 \\
\hline
\end{tabular}
\end{center}

\begin{center}
\small
\begin{tabular}{c c c c c c c c}
\hline
\(\lambda_{\mathrm{ridge}}\) & tr train & tr test & acc brute & F1 brute & acc oracle & F1 oracle & acc routée / ctx train / ctx test \\
\hline
0.2 & 15.468750000000911 & 15.479688630394707 & 0.7734375 & 0.7675757575757576 & 0.953125 & 0.9537112193362194 & 0.796875 / 0.859375 / 0.8515625 \\
0.8 & 20.687500000000085 & 20.5003307881078 & 0.7734375 & 0.7675757575757576 & 0.9765625 & 0.9765512265512266 & 0.8046875 / 0.859375 / 0.8515625 \\
1.5 & 21.952989309875143 & 21.937500000000078 & 0.7734375 & 0.7675757575757576 & 1.0 & 1.0 & 0.796875 / 0.8828125 / 0.8515625 \\
3.0 & 22.687499999998956 & 22.36718749999978 & 0.7734375 & 0.7675757575757576 & 1.0 & 1.0 & 0.78125 / 0.8828125 / 0.8515625 \\
\hline
\end{tabular}
\end{center}

Les gains exacts sont donc :
\[
\Delta_{\mathrm{oracle}} \in \{0.1796875,\ 0.203125,\ 0.2265625,\ 0.2265625\},
\]
\[
\Delta_{\mathrm{routed}} \in \{0.0234375,\ 0.03125,\ 0.0234375,\ 0.0078125\}.
\]

La monotonie énergétique reste nulle :
\[
\texttt{energy\_monotonic\_violations}=0
\]
dans tous les runs.

Les résidus de point fixe sont encore au niveau machine, entre
\[
5.719551864242296\times 10^{-17}
\quad\text{et}\quad
8.233436468072948\times 10^{-17}.
\]

Les traces de projecteur montent ici globalement d’environ \(15.5\) à \(22.7\), ce qui confirme que la géométrie active dépend de la campagne et de la jauge, même lorsque la loi scalaire normalisée est très stable.

\subsection{Ce que les chiffres disent exactement}

Trois faits ressortent de manière stable et doivent être conservés comme résultats centraux.

\paragraph{1. La loi normalisée scalaire est très fortement corroborée.}
Dans toutes les campagnes, pour toutes les valeurs de \(\lambda_{\mathrm{ridge}}\), le système vérifie
\[
\gamma=0.5,\qquad \tau=3.4,\qquad \rho_\star=\frac{1}{1+\tau}=0.22727272727272727.
\]
Les observations rapportées de \(\rho_{\mathrm{obs}}\) restent extrêmement proches de cette valeur ; les écarts sont petits, mais non uniformément nuls, et le fichier de résultats consigne précisément ces écarts par \(\texttt{rho\_gap}\), \(\texttt{operator\_gap}\), \(\texttt{projector\_gap}\) et \(\texttt{projector\_symmetry\_gap}\).

\paragraph{2. Le label \texttt{ridge} versus \texttt{none} est inactif numériquement dans les sorties collées.}
À \(\lambda_{\mathrm{ridge}}\) fixé, les résultats \texttt{ridge} et \texttt{none} coïncident à la précision reportée dans les trois campagnes.

\paragraph{3. Le routage reste la partie fragile du système.}
L’oracle PCH améliore fortement le probe et atteint même \(1.0\) dans plusieurs cas.
Le routage automatique, lui, reste en dessous de l’oracle dans toutes les campagnes, avec des gains modestes ou parfois nuls à faible \(\lambda_{\mathrm{ridge}}\).
Le bon contexte est donc la variable décisive.

\paragraph{4. La géométrie active varie avec \(\lambda_{\mathrm{ridge}}\) et avec la campagne.}
Les traces des projecteurs actifs et les gaps opérateur/projecteur ne sont pas constants d’une campagne à l’autre.  
Cela signifie que la loi scalaire est très stable, mais que la géométrie fine du hull actif n’est pas démontrée universelle par ces seules expériences.

\paragraph{5. La stabilité locale ne se confond pas avec la performance aval.}
La loi locale \(\rho_\star\) reste très stable, alors que les performances de classification varient encore selon le corpus, l’encodeur et la qualité du routage.  
Autrement dit, la validité locale du prox n’implique pas automatiquement la séparation aval.

\subsection{Interprétation théorique}

Les résultats soutiennent très fortement la version normalisée du Cas 3 :
\[
DT_{\eta,\alpha}^{\mathrm{PCH}}(u_\alpha^\star;w,c)
\approx
\bigl(I+\tau[I+2\gamma(I-P_\alpha)]\bigr)^{-1},
\qquad
\gamma=0.5,\ \tau=3.4.
\]

Autrement dit, la dynamique locale du noyau se comporte, dans la bonne jauge, comme un opérateur résolvant dont la partie scalaire est quasi invariante.

En revanche, les données fournies ne démontrent pas l’existence d’un projecteur limite universel \(P_0\) commun à toutes les campagnes, parce que :
\begin{itemize}
    \item les traces de projecteurs changent avec \(\lambda_{\mathrm{ridge}}\) ;
    \item les gaps opérateur/projecteur ne s’annulent pas uniformément ;
    \item les performances dépendent du split, de l’encodeur et du corpus.
\end{itemize}

Le bon énoncé scientifique est donc le suivant :
la normalisation rend la loi scalaire universelle à très haute précision sur les campagnes étudiées,
mais l’universalité géométrique complète de \(P_0\) reste ouverte.

Les limites expérimentales à garder explicitement sont :
\[
\text{un seul seed }(7),\qquad
\text{un seul split par campagne},\qquad
\text{un nombre fini de corpus et d’encodeurs}.
\]

\subsection{Conséquences conceptuelles, théoriques et pratiques}

\paragraph{Conséquence conceptuelle.}
PCH ne se comporte pas comme une heuristique libre : il se comporte comme un opérateur géométrique local fortement structuré.

\paragraph{Conséquence théorique.}
La partie la plus solide de la théorie est la loi résolvante normalisée
\[
\rho_\star=\frac{1}{1+\tau},
\]
validée numériquement sur les campagnes étudiées.
La partie encore ouverte est l’universalité complète du projecteur limite \(P_0\).

\paragraph{Conséquence pratique.}
Le vrai levier n’est pas d’augmenter \texttt{max\_iters} au-delà du régime où la convergence est déjà acquise.
Le levier principal est la qualité du routage de contexte.

\paragraph{Conséquence expérimentale.}
Le meilleur contexte améliore fortement le probe, parfois jusqu’à la séparation parfaite \(\texttt{accuracy}=1.0\), alors que le routage automatique demeure plus fragile et plus dépendant du split.

\subsection{Verdict final}

Le Cas 3 normalisé est \emph{très fortement corroboré} par les trois campagnes expérimentales.

La conclusion la plus robuste est :
\[
\boxed{
\text{la loi scalaire normalisée est très fortement corroborée,}
\quad
\text{mais l’universalité géométrique complète de } P_0 \text{ n’est pas démontrée.}
}
\]

En particulier :
\[
\boxed{
\text{Cas 3 normalisé : fortement corroboré expérimentalement comme loi résolvante normalisée robuste,}
\ \text{mais pas encore comme universalité géométrique absolue.}
}
\]




%%%%%%%%%%%%%%%%%%%%%%%%%%%%%%%%%%%%%%%%%%%%%%%%%%%%%%%%%%%%%%%%%%%%%%%%%%%%%%%%%%%%%%%%%%%%%%%%%%%%%%%%%%%%%%%%%%%%%%%%%%%%%%%%%%%%%%%%%%%%%%%%%%%%%%%%%%%%%%%%%%%%%%%%%%%%%%%%%%%%%%%%%%%%%%%%%%%%%%%%%%%%%%%%%%%%%%%%%%%%%%%%%%%%%%%%%%%%%%%%%%%%%%%%%%%%%%%%%%%%%%%%%%%%%%%%%%%%%%%%%%%demonstration cas 3 normalisé ++

\section{Réduction des hypothèses H2 et H3 à des axiomes plus fondamentaux}

\subsection{Point de départ : ce que fixent les expériences}

Les trois campagnes expérimentales montrent de manière robuste que la loi scalaire normalisée
\[
\rho_\star=\frac{1}{1+\tau}=\frac{1}{4.4}=0.22727272727272727
\]
est corroborée à très haute précision dans toutes les configurations testées, tandis que les quantités géométriques fines
\[
\operatorname{tr}(P_\alpha),\qquad \texttt{operator\_gap},\qquad \texttt{projector\_gap},\qquad \texttt{projector\_symmetry\_gap}
\]
restent variables selon le corpus, l'encodeur, la partition et \(\lambda_{\mathrm{ridge}}\).

Le bon découpage interprétatif est alors le suivant :
\begin{itemize}
    \item la partie scalaire normalisée du Cas 3 est stable ;
    \item la géométrie active fine n'est pas, à elle seule, démontrée universelle par les expériences ;
    \item H2 doit être lue comme une conséquence géométrique locale \emph{sous hypothèses de stabilité}, tandis que H3 exprime la concentration du projecteur ;
    \item la loi résolvante normalisée est l'objet directement validé par le protocole expérimental.
\end{itemize}

\subsection{Axiomes minimaux}

On remplace H2 et H3 par la structure minimale suivante.

\begin{assumption}[H0 : convexité et projection bien définie]
Pour chaque contexte \(c\), l'ensemble admissible \(M_\alpha\subset\mathbb{R}^{d_\alpha}\) est fermé, convexe et non vide.
La projection euclidienne \(\Pi_{M_\alpha}\) est donc bien définie et unique.
\end{assumption}

\begin{assumption}[H1 : régime régulier local]
Au point fixe pertinent \(u_\alpha^\star\), la projection \(\Pi_{M_\alpha}\) est différentiable dans un voisinage local suffisant, de sorte que
\[
\nabla^2 \dist(u_\alpha^\star,M_\alpha)^2
=
2\bigl(I-D\Pi_{M_\alpha}(u_\alpha^\star)\bigr).
\]
\end{assumption}

\begin{assumption}[H2' : concentration géométrique du hull]
Les embeddings de calibration d'un contexte donné sont modélisés par une famille aléatoire dont le hull empirique \(M_\alpha\) converge vers un hull population \(M_{\star,\alpha}\) en distance de Hausdorff, uniformément sur la famille considérée.
\end{assumption}

\begin{assumption}[H3\(_{\mathrm{geom}}\) : stabilité du projecteur limite géométrique]
Le projecteur différentiel du hull population converge vers un projecteur orthogonal limite déterministe \(P_0\), uniformément après normalisation :
\[
\|D\Pi_{M_{\star,\alpha}}(u_\alpha^\star)-P_0\|_{\mathrm{op}}\to 0.
\]
\end{assumption}

\begin{assumption}[H3\(_{\mathrm{res}}\) : stabilité du résolvant observable]
Le résolvant local normalisé
\[
R_\alpha:=\bigl(I+\tau[I+2\gamma(I-P_\alpha)]\bigr)^{-1}
\]
possède un observable scalaire \(s(R_\alpha)\), défini par le protocole expérimental comme la norme opératorielle du Jacobien proximal local,
\[
s(R_\alpha)=\|DT_{\eta,\alpha}(u_\alpha^\star;w,c)\|_{\mathrm{op}}.
\]
La valeur estimée par le protocole satisfait
\[
s(R_\alpha)\to \rho_\star=\frac{1}{1+\tau}.
\]
Dans le protocole expérimental, cet observable est approché par
\[
\rho_{\mathrm{obs}},\qquad \rho_{\mathrm{pred}},\qquad \rho_{\mathrm{gap}}.
\]
Par ailleurs, la comparaison matricielle complète est évaluée séparément par
\[
\texttt{operator\_gap}
=
\|DT_{\eta,\alpha}(u_\alpha^\star;w,c)-R_{\mathrm{pred}}\|_{\mathrm{op}}.
\]
\end{assumption}

\begin{remark}[Ce qui n'est pas le bon énoncé]
On ne prétend pas que le résolvant soit asymptotiquement un multiple scalaire de l'identité.
Au contraire, la cible opératorielle naturelle est
\[
R_\star:=\bigl(I+\tau[I+2\gamma(I-P_0)]\bigr)^{-1},
\]
qui est en général non scalaire, avec spectre
\[
\operatorname{Spec}(R_\star)
=
\left\{
\frac{1}{1+\tau}
\right\}^{r}
\cup
\left\{
\frac{1}{1+\tau(1+2\gamma)}
\right\}^{d_\alpha-r}.
\]
Le point expérimental validé n'est donc pas \(R_\alpha\approx \rho_\star I\), mais bien l'adhérence au résolvant normalisé correct, mesurée par l'observable scalaire et par la gap matricielle.
\end{remark}

\subsection{Définition des diagnostics numériques}

Le script estime par différences finies les jacobiens de la projection et du proximal, puis compare l'opérateur observé à l'opérateur prédit :
\[
R_{\mathrm{pred}}=\bigl(I+\tau[I+2\gamma(I-P_\alpha)]\bigr)^{-1}.
\]

Les diagnostics consignés dans les sorties sont :
\[
\rho_{\mathrm{obs}}=\|DT_{\eta,\alpha}(u_\alpha^\star;w,c)\|_{\mathrm{op}},
\qquad
\rho_{\mathrm{pred}}=\|R_{\mathrm{pred}}\|_{\mathrm{op}},
\qquad
\rho_{\mathrm{gap}}=|\rho_{\mathrm{obs}}-\rho_{\mathrm{pred}}|,
\]
\[
\texttt{operator\_gap}
=
\|DT_{\eta,\alpha}(u_\alpha^\star;w,c)-R_{\mathrm{pred}}\|_{\mathrm{op}},
\]
\[
\texttt{projector\_gap}
=
\|P_\alpha^2-P_\alpha\|_{\mathrm{op}},
\qquad
\texttt{projector\_symmetry\_gap}
=
\|P_\alpha-P_\alpha^\top\|_{\mathrm{op}},
\qquad
\texttt{projector\_trace}
=
\operatorname{tr}(P_\alpha).
\]

La loi scalaire \(\rho_\star=1/(1+\tau)\) est très fortement corroborée, mais les écarts opérateur/projecteur ne sont pas nuls uniformément sur toute la campagne ; ils sont précisément ceux que le fichier de résultats consigne.

\subsection{Proposition : H2 découle de H0, H1 et H2' sous une hypothèse supplémentaire de marge de stabilité}

\begin{proposition}[Stabilité géométrique des faces actives]
Supposons qu'il existe, pour le hull population \(M_{\star,\alpha}\), une face active \(F_{\star,\alpha}\) telle que le point projeté
\[
p_{\star,\alpha}:=\Pi_{M_{\star,\alpha}}(u_\alpha^\star)
\]
appartienne au relatif intérieur de \(F_{\star,\alpha}\), avec une marge géométrique positive :
\[
\operatorname{dist}\bigl(p_{\star,\alpha},\partial F_{\star,\alpha}\bigr)\ge \tau_0>0,
\]
et une séparation angulaire \(\alpha_0>0\) vis-à-vis des faces concurrentes.

Alors il existe une constante \(c_0>0\), dépendant seulement de \(\tau_0\) et \(\alpha_0\), telle que si
\[
d_H(M_\alpha,M_{\star,\alpha})\le c_0,
\]
alors la face active de \(M_\alpha\) au point \(u_\alpha^\star\) est combinatoirement stable, la projection \(\Pi_{M_\alpha}\) reste différentiable localement, et l'ensemble actif reste constant dans un voisinage de \(u_\alpha^\star\).

Autrement dit, H2 n'est pas un axiome indépendant : c'est une conséquence de la convexité, de la non-dégénérescence géométrique, de la concentration Hausdorff du hull, et d'une marge de stabilité des faces actives.
\end{proposition}

\begin{proof}
La marge \(\tau_0\) empêche le point projeté limite d'approcher le bord d'une face singulière, et la marge angulaire \(\alpha_0\) exclut la collision avec un cône normal concurrent.
Sous perturbation de Hausdorff petite devant ces marges, la face exposée par \(u_\alpha^\star\) reste la même à équivalence combinatoire près.

Comme la projection sur une face active fixe est localement régulière, sa différentiabilité locale est préservée et l'ensemble actif de contraintes ne change pas.
La stabilité de la face active est donc une conséquence standard de la non-dégénérescence et de la proximité Hausdorff.
\end{proof}

\subsection{Proposition : propagation de la convergence tangentielle}

\begin{proposition}[Concentration du différentiel de projection]
Sous les hypothèses H0, H1, H2' et H3\(_{\mathrm{geom}}\), on a
\[
\|D\Pi_{M_\alpha}(u_\alpha^\star)-P_0\|_{\mathrm{op}}\to 0
\]
uniformément sur la famille considérée.

En particulier, H3\(_{\mathrm{geom}}\) est la bonne hypothèse géométrique pour décrire la limite du projecteur différentiel.
\end{proposition}

\begin{proof}
D'après la proposition précédente, dès que \(d_H(M_\alpha,M_{\star,\alpha})\le c_0\), la face active de \(M_\alpha\) reste stable et
\[
D\Pi_{M_\alpha}(u_\alpha^\star)=P_{T_\alpha},
\]
où \(P_{T_\alpha}\) est le projecteur orthogonal sur l'espace tangent de la face active \(F_\alpha\).

Par H3\(_{\mathrm{geom}}\), on a convergence de ces projecteurs vers un projecteur limite déterministe \(P_0\). Donc
\[
\|D\Pi_{M_\alpha}(u_\alpha^\star)-P_0\|_{\mathrm{op}}
=
\|P_{T_\alpha}-P_0\|_{\mathrm{op}}
\to 0.
\]
\end{proof}

\subsection{Corollaire : le Hessien effectif devient un théorème}

Sous la normalisation
\[
\mu_\alpha=1+2\lambda_{\mathrm{ridge},\alpha},\qquad
\gamma_\alpha=\frac{\lambda_{\mathrm{ctx},\alpha}}{\mu_\alpha},\qquad
\tau_\alpha=\eta_\alpha\mu_\alpha,
\]
on obtient
\[
\widehat H_\alpha
=
I+2\gamma_\alpha(I-P_\alpha).
\]
Si \(\gamma_\alpha\to\gamma\) et \(P_\alpha\to P_0\), alors
\[
\widehat H_\alpha\to \widehat H_\star:=I+2\gamma(I-P_0).
\]
Par conséquent,
\[
DT_{\eta_\alpha,\alpha}(u_\alpha^\star;w,c)
=
\bigl(I+\tau_\alpha\widehat H_\alpha\bigr)^{-1}
\to
\bigl(I+\tau[I+2\gamma(I-P_0)]\bigr)^{-1}.
\]

Ainsi, le Hessien effectif universel n'est plus un axiome : il résulte de H0, H1, H2' et H3\(_{\mathrm{geom}}\), tandis que l'objet empirique directement corroboré est l'observable scalaire du résolvant, et, au niveau matriciel, la proximité du résolvant prédit \(R_{\mathrm{pred}}\).

\subsection{Lien exact avec les découvertes expérimentales}

Les expériences ex1, ex2 et ex3 montrent trois faits utiles pour la lecture théorique :

\begin{enumerate}
    \item La loi scalaire \(\rho_\star=1/(1+\tau)\) est corroborée à très haute précision dans toutes les campagnes.
    \item À \(\lambda_{\mathrm{ridge}}\) fixé, les variantes \texttt{ridge} et \texttt{none} sont numériquement identiques à la précision reportée, ce qui confirme que la bonne variable théorique est la jauge normalisée et non l'étiquette brute du regularizer.
    \item Les traces de projecteur et les gaps géométriques varient avec la campagne et avec \(\lambda_{\mathrm{ridge}}\), ce qui signifie que l'universalité absolue de \(P_0\) n'est pas encore établie expérimentalement. En revanche, la stabilité de l'opérateur résolvant prédit est précisément ce que les données rendent le plus robuste.
\end{enumerate}

La bonne conclusion est donc la suivante :

\begin{quote}
H2 peut être obtenue à partir d'une géométrie convexe minimale, d'une marge de stabilité des faces actives et d'un théorème de concentration des hulls empiriques.
H3\(_{\mathrm{geom}}\) reste la bonne hypothèse structurelle si l'on veut parler du projecteur limite \(P_0\).
Mais l'objet directement validé par les expériences est l'observable scalaire du résolvant, et, en parallèle, la proximité matricielle au résolvant prédit \(R_{\mathrm{pred}}\), pas la convergence brute de \(P_\alpha\) à elle seule.
\end{quote}

\subsection{Résumé quantitatif du canal scalaire}

Le diagnostic scalaire le plus important est l'écart à
\[
\rho_\star=\frac{1}{4.4}=0.22727272727272727.
\]
Sur les trois campagnes, les erreurs maximales observées sur \(\rho_{\mathrm{obs}}\) dans les valeurs agrégées de \texttt{normalization} valent :
\[
\max_{\mathrm{ex1}}|\rho_{\mathrm{obs}}-\rho_\star|
=
3.2508604025371923\times 10^{-6},
\]
\[
\max_{\mathrm{ex2}}|\rho_{\mathrm{obs}}-\rho_\star|
=
8.993255675715517\times 10^{-7},
\]
\[
\max_{\mathrm{ex3}}|\rho_{\mathrm{obs}}-\rho_\star|
=
1.4340065024320836\times 10^{-8}.
\]
Le maximum global sur l'ensemble des runs reportés est donc
\[
3.2508604025371923\times 10^{-6}.
\]
C'est cette stabilité scalaire qui est observée de manière la plus nette par le protocole.

\subsection{Traces de projecteur exactes}

Les traces de projecteur observées par le protocole, sur \emph{train} et sur \emph{test}, pour les quatre valeurs de \(\lambda_{\mathrm{ridge}}\), sont :

\begin{center}
\small
\begin{tabular}{c c c c c}
\hline
Campagne & \(\lambda_{\mathrm{ridge}}=0.2\) & \(\lambda_{\mathrm{ridge}}=0.8\) & \(\lambda_{\mathrm{ridge}}=1.5\) & \(\lambda_{\mathrm{ridge}}=3.0\) \\
\hline
ex1 train & 15.718750000000124 & 20.132812500000630 & 21.390624999999755 & 21.937569631649836 \\
ex1 test  & 15.521878804368042 & 20.312500000000350 & 21.523437499999620 & 21.929687499999694 \\
ex2 train & 15.167591529956862 & 20.132812500000630 & 21.390624999999755 & 21.937569631649836 \\
ex2 test  & 15.557274580859561 & 20.500330788107800 & 21.746580897633194 & 21.703197696365727 \\
ex3 train & 15.468750000000911 & 20.687500000000085 & 21.843749999999880 & 22.687499999998956 \\
ex3 test  & 15.479688630394707 & 20.500330788107800 & 21.746580897633194 & 22.367187499999780 \\
\hline
\end{tabular}
\end{center}

Ces valeurs rendent explicite le fait que la géométrie active varie avec la campagne et avec \(\lambda_{\mathrm{ridge}}\), sans remettre en cause la stabilité du canal scalaire. La variation n'est pas strictement monotone en \(\lambda_{\mathrm{ridge}}\) sur tous les sous-ensembles, mais l'amplitude des traces change nettement d'une campagne à l'autre.

\subsection{Version finale propre de la cible théorique}

La forme mathématiquement défendable du Cas 3 normalisé comporte désormais deux niveaux :

\[
s(R_\alpha)\to \rho_\star=\frac{1}{1+\tau}
\qquad\text{(observable scalaire directe)},
\]
et, au niveau opératoriel,
\[
R_\alpha \to R_\star:=\bigl(I+\tau[I+2\gamma(I-P_0)]\bigr)^{-1}
\qquad\text{(cible matricielle naturelle)}.
\]

Le bon schéma logique est donc :
\[
\text{convexité du hull}
\;\Longrightarrow\;
\text{projection bien définie}
\;\Longrightarrow\;
\text{marge géométrique + Hausdorff}
\;\Longrightarrow\;
\text{stabilité de la face active}
\;\Longrightarrow\;
D\Pi_{M_\alpha}(u_\alpha^\star)\to P_0
\;\Longrightarrow\;
H_{\mathrm{eff}}^\star
\;\Longrightarrow\;
DT_{\eta_\alpha,\alpha},
\]
avec, au niveau expérimental, un signal plus direct encore sur
\[
s(R_\alpha)\to \rho_\star
\]
et sur
\[
\|DT_{\eta,\alpha}(u_\alpha^\star;w,c)-R_{\mathrm{pred}}\|_{\mathrm{op}}.
\]

\subsection{Conclusion}

Dans la version normalisée, H2 est une conséquence analytique d'une géométrie convexe stable plus une concentration probabiliste des hulls empiriques, mais seulement sous une hypothèse explicite de marge de stabilité des faces actives.
H3\(_{\mathrm{geom}}\) reste l'hypothèse structurelle pour parler du projecteur limite \(P_0\), mais le phénomène le mieux corroboré par les expériences est l'observable scalaire du résolvant, avec en complément la proximité matricielle du résolvant normalisé prédit.
La théorie gagne ainsi en clarté : ce qui est fondamental, ce n'est pas l'universalité brute des hyperparamètres, mais la stabilité du résolvant associé ; la concentration géométrique vers \(P_0\) reste l'interprétation structurelle la plus naturelle, sans être encore le signal expérimental principal.





%%%%%%%%%%%%%%%%%%%%%%%%%%%%%%%%%%%%%%%%%%%%%%%%%%%%%%%%%%%%%%%%%%%%%%%%%%%%%%%%%%%%%%%%%%%%%%%%%%%%%%%%%%%%%%%%%%%%%%%%%%%%%%%%%%%%%%%%%%%%%%%%%%%%%%%%%%%%%%%%%%%%%%%%%%%%%%%%%%%%%%%%%%%%%%%%%%%%%%%%%%%%%%%%%%%%%%%%%%%%%%%%%%%%%%%%%%%%%%%%%%%%%%%%%%%%%%%%%%%%%%%%%%%%%%%%%%%%%%%%%%%





%%%%%%%%%%%%%%%%%%%%%%%%%%%%%%%%%%%%%%%%%%%%%%%%%%%%%%%%%%%%%%%%%%%%%%%%%%%%%%%%%%%%%%%%%%%%%%%%%%%%%%%%%%%%%%%%%%%%%%%%%%%%%%%%%%%%%%%%%%%%%%%%%%%%%%%%%%%%%%%%%%%%%%%%%%%%%%%%%%%%%%%%%%%%%%%%%%%%%%%%%%%%%%%%%%%%%%%%%%%%%%%%%%%%%%%%%%%%%%%%%%%%%%%%%%%%%%%%%%%%%%%%%%%%%%%%%%%%%%%%%%%





%%%%%%%%%%%%%%%%%%%%%%%%%%%%%%%%%%%%%%%%%%%%%%%%%%%%%%%%%%%%%%%%%%%%%%%%%%%%%%%%%%%%%%%%%%%%%%%%%%%%%%%%%%%%%%%%%%%%%%%%%%%%%%%%%%%%%%%%%%%%%%%%%%%%%%%%%%%%%%%%%%%%%%%%%%%%%%%%%%%%%%%%%%%%%%%%%%%%%%%%%%%%%%%%%%%%%%%%%%%%%%%%%%%%%%%%%%%%%%%%%%%%%%%%%%%%%%%%%%%%%%%%%%%%%%%%%%%%%%%%%%%





%%%%%%%%%%%%%%%%%%%%%%%%%%%%%%%%%%%%%%%%%%%%%%%%%%%%%%%%%%%%%%%%%%%%%%%%%%%%%%%%%%%%%%%%%%%%%%%%%%%%%%%%%%%%%%%%%%%%%%%%%%%%%%%%%%%%%%%%%%%%%%%%%%%%%%%%%%%%%%%%%%%%%%%%%%%%%%%%%%%%%%%%%%%%%%%%%%%%%%%%%%%%%%%%%%%%%%%%%%%%%%%%%%%%%%%%%%%%%%%%%%%%%%%%%%%%%%%%%%%%%%%%%%%%%%%%%%%%%%%%%%%





%%%%%%%%%%%%%%%%%%%%%%%%%%%%%%%%%%%%%%%%%%%%%%%%%%%%%%%%%%%%%%%%%%%%%%%%%%%%%%%%%%%%%%%%%%%%%%%%%%%%%%%%%%%%%%%%%%%%%%%%%%%%%%%%%%%%%%%%%%%%%%%%%%%%%%%%%%%%%%%%%%%%%%%%%%%%%%%%%%%%%%%%%%%%%%%%%%%%%%%%%%%%%%%%%%%%%%%%%%%%%%%%%%%%%%%%%%%%%%%%%%%%%%%%%%%%%%%%%%%%%%%%%%%%%%%%%%%%%%%%%%%





%%%%%%%%%%%%%%%%%%%%%%%%%%%%%%%%%%%%%%%%%%%%%%%%%%%%%%%%%%%%%%%%%%%%%%%%%%%%%%%%%%%%%%%%%%%%%%%%%%%%%%%%%%%%%%%%%%%%%%%%%%%%%%%%%%%%%%%%%%%%%%%%%%%%%%%%%%%%%%%%%%%%%%%%%%%%%%%%%%%%%%%%%%%%%%%%%%%%%%%%%%%%%%%%%%%%%%%%%%%%%%%%%%%%%%%%%%%%%%%%%%%%%%%%%%%%%%%%%%%%%%%%%%%%%%%%%%%%%%%%%%%





%%%%%%%%%%%%%%%%%%%%%%%%%%%%%%%%%%%%%%%%%%%%%%%%%%%%%%%%%%%%%%%%%%%%%%%%%%%%%%%%%%%%%%%%%%%%%%%%%%%%%%%%%%%%%%%%%%%%%%%%%%%%%%%%%%%%%%%%%%%%%%%%%%%%%%%%%%%%%%%%%%%%%%%%%%%%%%%%%%%%%%%%%%%%%%%%%%%%%%%%%%%%%%%%%%%%%%%%%%%%%%%%%%%%%%%%%%%%%%%%%%%%%%%%%%%%%%%%%%%%%%%%%%%%%%%%%%%%%%%%%%%





%%%%%%%%%%%%%%%%%%%%%%%%%%%%%%%%%%%%%%%%%%%%%%%%%%%%%%%%%%%%%%%%%%%%%%%%%%%%%%%%%%%%%%%%%%%%%%%%%%%%%%%%%%%%%%%%%%%%%%%%%%%%%%%%%%%%%%%%%%%%%%%%%%%%%%%%%%%%%%%%%%%%%%%%%%%%%%%%%%%%%%%%%%%%%%%%%%%%%%%%%%%%%%%%%%%%%%%%%%%%%%%%%%%%%%%%%%%%%%%%%%%%%%%%%%%%%%%%%%%%%%%%%%%%%%%%%%%%%%%%%%%





%%%%%%%%%%%%%%%%%%%%%%%%%%%%%%%%%%%%%%%%%%%%%%%%%%%%%%%%%%%%%%%%%%%%%%%%%%%%%%%%%%%%%%%%%%%%%%%%%%%%%%%%%%%%%%%%%%%%%%%%%%%%%%%%%%%%%%%%%%%%%%%%%%%%%%%%%%%%%%%%%%%%%%%%%%%%%%%%%%%%%%%%%%%%%%%%%%%%%%%%%%%%%%%%%%%%%%%%%%%%%%%%%%%%%%%%%%%%%%%%%%%%%%%%%%%%%%%%%%%%%%%%%%%%%%%%%%%%%%%%%%%





%%%%%%%%%%%%%%%%%%%%%%%%%%%%%%%%%%%%%%%%%%%%%%%%%%%%%%%%%%%%%%%%%%%%%%%%%%%%%%%%%%%%%%%%%%%%%%%%%%%%%%%%%%%%%%%%%%%%%%%%%%%%%%%%%%%%%%%%%%%%%%%%%%%%%%%%%%%%%%%%%%%%%%%%%%%%%%%%%%%%%%%%%%%%%%%%%%%%%%%%%%%%%%%%%%%%%%%%%%%%%%%%%%%%%%%%%%%%%%%%%%%%%%%%%%%%%%%%%%%%%%%%%%%%%%%%%%%%%%%%%%%





%%%%%%%%%%%%%%%%%%%%%%%%%%%%%%%%%%%%%%%%%%%%%%%%%%%%%%%%%%%%%%%%%%%%%%%%%%%%%%%%%%%%%%%%%%%%%%%%%%%%%%%%%%%%%%%%%%%%%%%%%%%%%%%%%%%%%%%%%%%%%%%%%%%%%%%%%%%%%%%%%%%%%%%%%%%%%%%%%%%%%%%%%%%%%%%%%%%%%%%%%%%%%%%%%%%%%%%%%%%%%%%%%%%%%%%%%%%%%%%%%%%%%%%%%%%%%%%%%%%%%%%%%%%%%%%%%%%%%%%%%%%
























































































%%%%%%%%%%%%%%%%%%%%%%%%%%%%%%%%%%%%%%%%%%%%%%%%%%%%%%%%%%%%%%%%%%%%%%%%%%%%%%%%%%%%%%%%%%%%%%%%%%%%%%%%%%%%%%%%%%%%%%%%%%%%%%%%%%%%%%%%%%%%%%%%%%%%%%%%%%%%%%%%%%%%%%%%%%%%%%%%%%%%%%%%%%%%%%%%%%%%%%%%%%%%%%%%%%%%%%%%%%%%%%%%%%%%%%%%%%%%%%%%%%%%%%%%%%%%%%%%%%%%%%%%%%%%%%%%%%%%%%%%%%%confrontation avec la theorie complete :

\section{Évaluation du niveau scientifique de la théorie complète et confrontation aux résultats expérimentaux}

\subsection{Préambule méthodologique}

L'évaluation présentée dans cette section porte sur la théorie complète décrite dans le document fondateur, et non uniquement sur le noyau PCH.

Une première lecture centrée exclusivement sur les expériences pourrait conduire à identifier PCH comme l'objet principal de la théorie. Une analyse intégrale du document montre cependant que le noyau PCH n'est qu'un composant d'une architecture théorique beaucoup plus vaste.

La théorie développée dans le document est une théorie \emph{PCH--ESM} complète comprenant plusieurs niveaux hiérarchiques :

\begin{enumerate}
    \item la théorie PH ;
    \item la théorie PCH ;
    \item la décomposition spectrale en espaces spécialisés ;
    \item le mécanisme de collage global \(G\) ;
    \item la théorie Multi-PH ;
    \item la modélisation hiérarchique de la croyance, de l'incertitude et de la représentation du monde.
\end{enumerate}

Par conséquent, la question scientifique correcte n'est pas :

\begin{quote}
``Les expériences valident-elles PCH ?''
\end{quote}

mais plutôt :

\begin{quote}
``Quelle partie de la théorie complète est démontrée analytiquement, quelle partie est soutenue expérimentalement, et quelle partie reste ouverte ?''
\end{quote}

\subsection{Valeur scientifique de la théorie seule}

On considère ici uniquement le contenu mathématique du document, indépendamment des expériences réalisées.

\subsubsection{Bloc PH}

La théorie PH constitue l'une des parties les plus solides de l'édifice.

Les résultats démontrés comprennent notamment :

\begin{itemize}
    \item existence du minimiseur ;
    \item unicité ;
    \item stabilité ;
    \item convergence proximale ;
    \item propriétés de contractivité ;
    \item existence des points fixes.
\end{itemize}

Ces résultats reposent essentiellement sur des outils standards d'analyse convexe, de théorie des opérateurs monotones et de théorie proximale.

Du point de vue mathématique, cette partie peut être considérée comme rigoureuse et mature.

\subsubsection{Bloc PCH}

Le noyau PCH hérite de la structure analytique de PH tout en introduisant une composante géométrique liée aux hulls convexes.

Le document établit notamment :

\begin{itemize}
    \item la purification contextuelle ;
    \item la définition variationnelle du noyau ;
    \item l'existence du minimiseur ;
    \item son unicité ;
    \item la stabilité de la dynamique associée ;
    \item les propriétés de convergence du schéma proximal.
\end{itemize}

Cette partie est également très solide du point de vue mathématique.

\subsubsection{Bloc de spécialisation spectrale}

Le document introduit une décomposition spectrale de la forme

\[
\mathcal H
=
\mathcal H_0
\oplus
\mathcal H_1
\oplus
\cdots
\oplus
\mathcal H_n.
\]

Cette construction repose sur des projecteurs spectraux et sur une décomposition orthogonale de l'espace latent.

Les résultats de cohérence mathématique apparaissent rigoureux.

Toutefois, contrairement à PH et PCH, cette partie dispose actuellement de peu de validations empiriques directes.

\subsubsection{Bloc de collage global \(G\)}

Le document introduit un mécanisme de reconstruction globale à partir des composantes spécialisées.

Les propriétés démontrées comprennent notamment :

\begin{itemize}
    \item cohérence du collage ;
    \item existence du minimiseur global ;
    \item unicité ;
    \item optimalité de type inverse-variance ;
    \item stabilité de la reconstruction.
\end{itemize}

Mathématiquement, cette partie apparaît cohérente et relativement développée.

\subsubsection{Bloc Multi-PH}

Le niveau Multi-PH constitue l'ambition la plus élevée du document.

Il introduit notamment :

\begin{itemize}
    \item des représentations hiérarchiques du monde ;
    \item des mécanismes de croyance ;
    \item des mécanismes d'incertitude ;
    \item des structures de composition entre espaces spécialisés ;
    \item une dynamique hiérarchique globale.
\end{itemize}

La formalisation mathématique est cohérente.

En revanche, cette partie ne bénéficie pas encore du même niveau de validation que PH ou PCH.

\subsection{Évaluation de la théorie seule}

Au vu des démonstrations disponibles, la hiérarchie suivante paraît raisonnable :

\begin{center}
\begin{tabular}{lc}
\hline
Bloc théorique & Niveau de maturité \\
\hline
PH & très élevé \\
PCH & très élevé \\
Spécialisation spectrale & élevé \\
Collage \(G\) & élevé \\
Multi-PH & moyen à élevé \\
Cas 3 universel & ouvert \\
\hline
\end{tabular}
\end{center}

\subsection{Apport des expériences réalisées}

Les campagnes expérimentales apportent une information nouvelle qui n'était pas contenue dans le document théorique seul.

Elles permettent une confrontation directe entre :

\begin{itemize}
    \item les prédictions analytiques ;
    \item les comportements numériques observés.
\end{itemize}

\subsection{Statut du Cas 1}

Les expériences montrent :

\begin{itemize}
    \item absence de violation énergétique ;
    \item convergence systématique ;
    \item résidus de point fixe compris entre \(10^{-17}\) et \(10^{-16}\).
\end{itemize}

Le Cas 1 est donc simultanément :

\[
\boxed{
\text{démontré analytiquement}
}
\]

et

\[
\boxed{
\text{validé expérimentalement}
}
\]

avec un niveau de précision exceptionnel.

\subsection{Statut du Cas 2}

La théorie prédit une réduction géométrique de la distance au hull admissible.

Les expériences montrent notamment :

\[
0.1532
\rightarrow
0.00746
\]

dans le cas oracle de l'expérience ex1.

La réduction est :

\begin{itemize}
    \item systématique ;
    \item forte ;
    \item stable ;
    \item reproductible.
\end{itemize}

Le Cas 2 doit donc être considéré comme

\[
\boxed{
\text{très fortement supporté expérimentalement}
}
\]

bien qu'il ne dispose pas encore d'une preuve universelle complète.

\subsection{Statut du Cas 3 local}

La découverte principale issue des expériences concerne la loi locale :

\[
\rho(\eta)
=
\frac{1}{1+\eta\mu_{\mathrm{eff}}}.
\]

Cette relation est observée avec des erreurs relatives de l'ordre de

\[
10^{-16}.
\]

Une telle précision est extrêmement rare dans un système réel impliquant :

\begin{itemize}
    \item des embeddings textuels ;
    \item des hulls convexes ;
    \item des solveurs numériques ;
    \item des données réelles.
\end{itemize}

Par conséquent :

\[
\boxed{
\text{Cas 3 local très fortement supporté}
}
\]

et probablement la découverte expérimentale la plus importante de l'ensemble du programme.

\subsection{Ce qui est réellement validé aujourd'hui}

Les expériences réalisées testent essentiellement :

\[
\boxed{
\text{PH}
\longrightarrow
\text{PCH}
}
\]

et très peu les autres composantes du système théorique.

En particulier, les expériences ne testent pratiquement pas :

\begin{itemize}
    \item la spécialisation spectrale complète ;
    \item la partition canonique ;
    \item le collage optimal \(G\) ;
    \item les compromis capacité/spécialisation ;
    \item Multi-PH ;
    \item la théorie hiérarchique de la représentation du monde.
\end{itemize}

Il serait donc scientifiquement incorrect d'affirmer aujourd'hui que la théorie complète est validée expérimentalement.

Ce qui est validé est essentiellement :

\[
\boxed{
\text{le bloc PH/PCH}
}
\]

et certaines conséquences locales associées.

\subsection{Évaluation globale}

\subsubsection{Théorie seule}

La théorie prise isolément présente :

\begin{itemize}
    \item une structure mathématique cohérente ;
    \item des définitions précises ;
    \item des théorèmes nombreux ;
    \item une architecture conceptuelle ambitieuse ;
    \item une absence de contradiction interne manifeste.
\end{itemize}

Son niveau scientifique peut être considéré comme élevé.

\subsubsection{Théorie plus expériences}

Les expériences apportent :

\begin{itemize}
    \item la validation du Cas 1 ;
    \item un fort soutien au Cas 2 ;
    \item une validation numérique extrêmement précise du Cas 3 local ;
    \item une caractérisation empirique de
    \[
    \mu_{\mathrm{eff}};
    \]
    \item l'identification du routage comme principal facteur limitant.
\end{itemize}

L'ensemble théorie + expériences possède donc une valeur scientifique significativement supérieure à celle du document seul.

\subsection{Résultat le plus important obtenu à ce jour}

Le résultat le plus important n'est probablement ni :

\begin{quote}
``PCH converge''
\end{quote}

ni

\begin{quote}
``PCH améliore l'accuracy''
\end{quote}

mais plutôt l'observation suivante :

\[
\boxed{
DT_\eta(u^\star)
\approx
(I+\eta H_{\mathrm{eff}})^{-1}
}
\]

avec une précision numérique proche de la précision machine.

Cette observation relie directement :

\begin{enumerate}
    \item la théorie variationnelle ;
    \item la géométrie des hulls convexes ;
    \item la dynamique locale du proximal ;
    \item les résultats observés sur des données réelles.
\end{enumerate}

Elle constitue actuellement le lien le plus fort entre la théorie abstraite et les observations expérimentales.

\subsection{Verdict final}

À l'heure actuelle, la situation scientifique peut être résumée ainsi :

\[
\boxed{
\text{PH : démontré et validé}
}
\]

\[
\boxed{
\text{PCH : démontré et validé}
}
\]

\[
\boxed{
\text{Cas 1 : confirmé}
}
\]

\[
\boxed{
\text{Cas 2 : fortement supporté}
}
\]

\[
\boxed{
\text{Cas 3 local : très fortement supporté}
}
\]

\[
\boxed{
\text{Cas 3 universel : ouvert}
}
\]

\[
\boxed{
\text{Spécialisation spectrale : peu testée}
}
\]

\[
\boxed{
\text{Collage }G\text{ : peu testé}
}
\]

\[
\boxed{
\text{Multi-PH : essentiellement non testé}
}
\]

La conclusion correcte n'est donc pas que la théorie complète est démontrée.

La conclusion correcte est que le bloc PH/PCH est désormais fortement étayé à la fois analytiquement et expérimentalement, tandis que les niveaux supérieurs de la théorie PCH--ESM restent encore largement à explorer et à confronter aux données.